% concatenated from template.tex
\documentclass{article}



\usepackage{arxiv}

\usepackage[utf8]{inputenc} % allow utf-8 input
\usepackage[T1]{fontenc}    % use 8-bit T1 fonts
\usepackage[hidelinks]{hyperref}       % hyperlinks
\usepackage{url}            % simple URL typesetting
\usepackage{booktabs}       % professional-quality tables
\usepackage{amsfonts}       % blackboard math symbols
\usepackage{nicefrac}       % compact symbols for 1/2, etc.
\usepackage{microtype}      % microtypography
\usepackage{lipsum}		% Can be removed after putting your text content
\usepackage{graphicx}
\usepackage{natbib}
\usepackage{doi}

\usepackage{amsfonts}
\usepackage{stmaryrd}
\usepackage{amsmath,amssymb}
\usepackage{mathtools}
\usepackage{amsthm}
% \usepackage{natbib}
% \usepackage[square,numbers]{natbib}
\usepackage{enumitem}   
\usepackage{cancel}
\usepackage{subcaption}
\usepackage{float} 


% \DeclareMathOperator*{\argmax}{argmax}
\usepackage{xcolor}
\newtheorem{theorem}{Theorem}
\newtheorem{lemma}{Lemma}
\newtheorem{corollary}{Corollary}
\newtheorem{proposition}{Proposition}

\theoremstyle{definition}
\newtheorem{definition}{Definition}[section]
\newtheorem{conjecture}{Conjecture}[section]
\newtheorem{example}{Example}[section]
% \newtheorem{theorem}{Theorem}[section]
% \newtheorem{lemma}{Lemma}[section]
% \newtheorem{proposition}{Proposition}[section]
% \newtheorem{corollary}{Corollary}[section]
\newtheorem{aim}{Aim}[section]
\newtheorem{question}{Question}[section]
\newtheorem{task}{Task}
\newtheorem{observation}{Observation}[section]
\newtheorem{remark}{Remark}[section]


\newcommand{\pr}[1]{ \left( #1 \right) }
\newcommand{\inner}[2]{\langle #1, #2 \rangle}
\newcommand{\derivative}[1]{#1^{\prime}}
\newcommand{\abs}[1]{\left| #1 \right|} 
\newcommand{\set}[1]{\left\{#1\right\}}
\newcommand{\ie}{i.e.}

\newcommand{\js}[1]{{\color{purple}#1}}


\title{Numerical Analysis of HiPPO-LegS ODE for\\
Deep State Space Models
}

%\date{September 9, 1985}	% Here you can change the date presented in the paper title
%\date{} 					% Or removing it

% \author{ {
% % \includegraphics[scale=0.06]{orcid.pdf}\hspace{1mm}
% Jaesung R. Park}
% % \thanks{Use footnote for providing further
% % 		information about author (webpage, alternative
% % 		address)---\emph{not} for acknowledging funding agencies.} 
%     \\
% 	% Department of
%     Mathematical Sciences\\
% 	Seoul National University\\
% 	% Pittsburgh, PA 15213 \\
% 	\texttt{ryanpark7@snu.ac.kr} \\
% 	%% examples of more authors
%     % \hspace{-7.3in}
% 	\And
%     % \hspace{-7.3in}
%     {
%     % \includegraphics[scale=0.06]{orcid.pdf}\hspace{1mm}
%     Jaewook J. Suh} \\
% 	% Department of 
%      Computational Applied Mathematics \& Operations Research \\
%     % CMOR\\
% 	Rice University\\
% 	\texttt{jacksuh@rice.edu} \\
%     \And
%     {
%     % \includegraphics[scale=0.06]{orcid.pdf}    \hspace{1mm}
%     Ernest K. Ryu} \\
% 	% Department of 
%     Mathematics\\
% 	University of California, Los Angeles\\
% 	\texttt{eryu@math.ucla.edu} \\
% 	%% \AND
% 	%% Coauthor \\
% 	%% Affiliation \\
% 	%% Address \\
% 	%% \texttt{email} \\
% }
\author{
\normalfont
% \hspace{-0.4in}
\begin{tabular}{@{}c@{\hspace{1em}}c@{\hspace{0.7em}}c@{\hspace{0.2em}}c@{}}  % @{} removes default spaces
    \textbf{Jaesung R. Park} & \textbf{Jaewook J. Suh} & \textbf{Youngjoon Hong} &\textbf{Ernest K. Ryu} \\
    Mathematical Sciences &
    CMOR & Mathematical Sciences & Mathematics \\
    Seoul National University & Rice University& Seoul National University & UCLA \\
    \texttt{ryanpark7@snu.ac.kr} & \texttt{jacksuh@rice.edu} &\texttt{hongyj@snu.ac.kr}& \texttt{eryu@math.ucla.edu}
\end{tabular}
}

% Uncomment to remove the date
\date{}

% Uncomment to override  the `A preprint' in the header
\renewcommand{\headeright}{ }
%\renewcommand{\undertitle}{Technical Report}
\renewcommand{\shorttitle}{ }

%%% Add PDF metadata to help others organize their library
%%% Once the PDF is generated, you can check the metadata with
%%% $ pdfinfo template.pdf
\hypersetup{
pdftitle={LegS arXiv},
pdfsubject={ },
pdfauthor={Jaesung R. Park, Jaewook J. Suh, Ernest K. Ryu},
pdfkeywords={ },
}

\begin{document}
\maketitle

	% The HiPPO framework is a method for an online approximation of sequential data, which has empirically experienced lots of success in enhancing the LSTM structure, and is now the backbone of many state-of-the-art state space models, including S4 and mamba. While the HiPPO method allows various measures into its structure, the Leg-S (Legendre Scaled) formulation has distinct advantages with less hyperparameter tuning and scale invariance. However, the uniform measure scaled by $1/t$ introduces a singularity at $t=0$ that is not found in other methods. In this work, we have mathematically verified the well-posedness of the LegS structure and performed an error analysis through various discretization methods. We verify that the careful consideration of the initial condition is crucial for handling the singularity both in theoretical and empirical terms. Moreover, we prove that the discretization error could be controlled despite its singularity.
    
\begin{abstract}
In deep learning, the recently introduced state space models utilize HiPPO (High-order Polynomial Projection Operators) memory units to approximate continuous-time trajectories of input functions using ordinary differential equations (ODEs), and these techniques have shown empirical success in capturing long-range dependencies in long input sequences. However, the mathematical foundations of these ODEs, particularly the singular HiPPO-LegS (Legendre Scaled) ODE, and their corresponding numerical discretizations remain unsettled. In this work, we fill this gap by establishing that HiPPO-LegS ODE is well-posed despite its singularity, albeit without the freedom of arbitrary initial conditions. Further, we establish convergence of the associated numerical discretization schemes for Riemann integrable input functions.
\end{abstract}







% keywords can be removed
% \keywords{HiPPO-LegS \and Well-posedness of ODE \and Discretization convergence analysis}


\section{Introduction} \label{sec:intro}
% Specifically, a linear state-space representation has the form
% \begin{equation}
%     \begin{aligned} \label{eq:ss_representation}
%     \dot{x}(t) &= A(t)x(t) + B(t)f(t)\\
%     y(t) &= C(t)x(t) + D(t)f(t) %\nonumber
% \end{aligned}
% \end{equation}

% Here, $x(t) \in \mathbb{R}^{N}$ is the state vector, $y(t) \in \mathbb{R}^{p}$ is the output, and $f(t) \in \mathbb{R}^q$ is the input function.
% The matrices $A(t) \in \mathbb{R}^{N \times N}, B(t) \in \mathbb{R}^{q \times N}, C(t) \in \mathbb{R}^{N \times p}, D(t) \in \mathbb{R}^{q \times p}$ encode how the system evolves and responds to inputs.
% When all the matrices $A, B, C, D$ do not explicitly depend on time, \eqref{eq:ss_representation} is called a linear time-invariant (LTI) system.


%ODE 문제일지라도 왜 중요한지. 정보를 확실하게.
%potential downstream impact

%Toy model로 실험해서 log-log plot에서 실험과 증명과 일치한다


% State-space representation is a classical mathematical framework widely used to model physical systems in diverse fields such as control engineering, signal processing, and computational neuroscience, and there has been much recent interest in using state-space representations for sequence modeling in deep learning \citep{gu_efficiently_2022, dao2024transformers, zhu2024vision, nguyen2022s4nd, goel_its_2022}.
% The foundation of these models could be understood as the combination of classical state space representation \citep{williams2007linear, zak2003systems} and the HiPPO (high-order polynomial projection operators) theory \citep{gu2020hippo}.
State-space representation is a cornerstone of dynamical-system theory and has been instrumental in the analysis and control of physical processes in control engineering, signal processing, and computational neuroscience. In the deep-learning literature, this classical framework has recently re-emerged as a promising paradigm for sequence modelling, offering a principled alternative to recurrent and attention-based architectures \citep{gu_efficiently_2022, dao2024transformers, zhu2024vision, nguyen2022s4nd, goel_its_2022}.
Modern state-space models for long sequences build on a synthesis of two pillars: (i) linear state-space theory in its canonical form \citep{williams2007linear, zak2003systems} and (ii) the HiPPO (High-order Polynomial Projection) framework \citep{gu2020hippo}, which prescribes optimal polynomial projections for compressing the history of an input signal. This amalgamation provides both an interpretable memory mechanism and a mathematically tractable route for capturing long-range dependencies within deep architectures.

HiPPO is a framework using an $N$-dimensional ordinary differential equation (ODE) to approximate the continuous-time history of an input function $f$.
% By keeping the $N$ projections of the input functions with respect to the polynomial basis, the model could efficiently perform an online approximation of the input function.
In particular, the HiPPO-LegS (Legendre Scaled) ODE is
\[
\derivative{c}(t) = -\frac{1}{t}Ac(t) + \frac{1}{t}Bf(t),
\] 
for $t\in [0,T]$, where $T>0$ is some terminal time and $f \colon  [0,T] \to \mathbb{R}$ is an input function.
With a specific choice of $A \in \mathbb R^{N \times N}$ and $B \in \mathbb R^{N \times 1}$, the solution $c\colon [0,T]\rightarrow \in \mathbb{R}^N$ encodes the continuous-time history of $f$ via
\begin{align*} 
    c_j(t) &= \frac{1}{t} \Big\langle f(\cdot ), \sqrt{2j-1} P_{j-1}\big(\tfrac{2\,\cdot} {t}-1\big) \Big\rangle_{L^2([0,t])}
    =
    \frac{\sqrt{2j-1}}{t} \int_0^t f(s) P_{j-1}\big(\tfrac{2s} {t}-1\big) \;ds,
\end{align*}
where $P_{j-1}$ is the $j-1$-th Legendre polynomial, for $j=1,2,\dots$.
% \citep{legendre1785recherches, szeg1939orthogonal}.
% {\color{red}[ XXX provide citationXXX ]}.
The input function $f$ can be approximately reconstructed by the formula
\begin{align*}
    f(s) \approx \sum_{j=1}^Nc_j(t) \sqrt{2j-1} P_{j-1}\big(\tfrac{2s}{t}-1\big)\qquad\text{ for }s\in [0,t].
\end{align*}
Since $\{ \sqrt{2n+1} P_n(\frac{2\,\cdot}{t}-1)\}_{n=0}^\infty$ forms an
% {\color{red} XXX orthonormal? XXX}
orthonormal polynomial basis on the $\frac{1}{t} \mathbb{I}_{[0,t]}(s)ds$ measure, this approximation can be viewed as the projection of $f$ on to the $N$-dimensional subspace spanned by $\{\sqrt{2n+1} P_n(\frac{2\,\cdot}{t}-1)\}_{n=0}^{N-1}$.

% \eqref{eq:LegS_projection} could be understood as an orthogonal projection of $f \colon  [0, t] \to \mathbb{R}$. 

 % The HiPPO theory generalized this framework by utilizing various measures and orthogonal polynomials that the input function is projected to. In particular, the HiPPO-LegS(Legendre Scaled) formulation used a measure with time-varying length, with the state-space representation as below:

% where $c \in \mathbb{R}^{N}$ is the state vector and $A \in \mathbb{R}^{N \times N}, B \in \mathbb{R}^{N \times 1}$ are pre-determined matrices. $f(t) \in \mathbb{R}$ is the scalar input function of which the model aims to approximate its trajectory. 
Unlike previous linear time-invariant (LTI) methods such as Legendre Memory Units (LMUs) \citep{voelker2019legendre}, which considers measures with a fixed-length support, the LegS formulation considers the uniform measure on $[0,t]$ that widens with the progression of time $t$, and therefore provides a memory unit keeping track of the entire trajectory of $f(\cdot)$ from time $0$ to $t$.
% formulation projects the input function $f(t)$ onto the basis $\mathbb{I}_{[0,t]}$, 
% allows the model to keep track of the whole trajectory of $f(t)$.
This distinctive property makes the LegS formulation powerful in many practical applications.

However, despite receiving much attention for its use in state space models in deep learning, the careful mathematical foundation of the LegS ODE is missing. To begin with, the singularity at $t=0$ renders the question of existence and uniqueness of the solution $c(t)$ a non-obvious matter. Moreover, it is unclear if the numerical methods used in the work of \cite{gu2020hippo} are mathematically justified in the sense of convergence: In the limit of small stepsizes, do the discrete simulations converge to the true continuous-time solution? What regularity conditions must $f$ satisfy for such convergence?


\paragraph{Contributions.} 
In this work, we provide the rigorous mathematical foundations of the HiPPO-LegS ODE formulation and its discretization. Specifically, we show that (i) the solution to the LegS ODE exists and is unique, but the initial condition is fixed to a predetermined value depending on $f(0)$ and (ii) the commonly used discretization schemes for LegS converges to the exact continuous-time solution for all Riemann integrable $f$.


% For our second contribution, our proof technique is to show that the numerical integration schemes reduces to a quadrature of $f$, we establish a general framework for understanding the convergence of discretization methods in this context.


% , the integrator of the ODE is not Lipschitz continuous with respect to $c(t)$, and hence we cannot assume the existence and the uniqueness of the solution \textit{a priori}. Moreover, when used in practice, the state space representation \eqref{eq:intro_legs} has to be discretized to be processed in a computer. This could also be problematic in the sense that the discrete-time version state space model may not converge to the continuous-time state space model without additional assumptions.

% to construct a computationally efficient and generally applicable framework for generative modeling.\citep{gu2023modeling, gu2021combining}


\subsection{Related works} \label{sec:related_works}

\paragraph{State space models for deep learning.}
The use of state space models (SSMs) in deep learning has gained significant recent attention due to their ability to process sequential data efficiently.
While the transformers architecture \citep{vaswani2017attention} has become the standard for language models, recent SSM models such as mamba  \citep{gu_mamba_2023} have been reported to achieve state-of-the-art results, especially in handling long sequences.

Large-scale SSMs deploy an initialization scheme motivated by the HiPPO theory \citep{gu2023train}. One distinctive characteristic of state-of-the-art SSMs is that the computation cost displays a near-linear growth with respect to sequence length, unlike the quadratic growth of transformers. S4 \citep{gu_efficiently_2022} uses the fast Fourier transform to attain the near-linear cost, whereas Mamba leverages hardware-aware computation techniques to attain near-linear parallel compute steps. The SSM architecture has been applied to or motivated numerous model structures \citep{fu2022hungry, hasani2022liquid, sun2023retentive, peng2023rwkv} and are used across various modalities \citep{zhu2024vision, li2025videomamba, shams2024ssamba}. 


\paragraph{Legendre memory units for LSTMs.}
A fundamental challenge in training recurrent neural networks (RNNs) is the vanishing gradient problem, which causes long-range dependencies in temporal data to be lost during training  \citep{bengio1994learning, le2015simple}. While LSTMs \citep{schmidhuber1997long} alleviate this problem by incorporating nonlinear gating mechanisms, modeling very long sequences remains challenging.
Motivated by applications in computational neuroscience, LMUs \citep{voelker2019dynamical, voelker2019legendre} introduced a novel approach to extend LSTM's capability to `remember' the sequence information by using a $N$-dimensional ODE. Using Pad\'{e} approximants \citep{pade1892representation}, the ODE is constructed so that the solution is the projection of the input function on the orthonormal basis of measure $\mathbb{I}_{[t-\theta, t]}$, where $\theta$ is a hyperparameter. 
The HiPPO framework could be understood as a generalization of LMUs. 
While LMU and its variants has proven to be effective for long sequence modeling \citep{liu_lmuformer_2024,zhang2023performance,  chilkuri_parallelizing_2021}, their scope is limited to LTI methods.


% well. Specifically, they derived a state space model represented by the following ODE
% \begin{align} \label{eq:LMU_ODE}
%     \theta \dot{c}(t) &= Ac(t) + Bf(t)
% \end{align}
% where $c(t) \in \mathbb{R}^N$ as the state vector and $f(t) \in \mathbb{R}$ the input function. $A \in \mathbb{R}^{N \times N}$ and $B \in \mathbb{R}^{N \times 1}$ are predetermined matrices derived by the use of Pad\'{e} approximants\citep{pade1892representation}.
% A key feature of this ODE is that the solution $c(t)$ becomes the projection of $f(t)$ on the domain $[t, t-\theta]$. Specifically, the $j$-th component of $c$ becomes
% \begin{align} \label{eq:LMU_projection}
%     c(t)_j &= \Big\langle f(t-\theta'), P_{j-1}\pr{\frac{2\theta'}{\theta}-1} \Big\rangle,
%     \qquad 0 \leq \theta' \leq \theta
% \end{align}
% where $P_j$ is the $j$-th Legendre polynomial. Since $\{P_n\pr{\frac{2\theta'}{\theta}-1}\}_{n \in \mathbb{N}}$ form an orthogonal basis on $[0, \theta]$, \eqref{eq:LMU_projection} could be understood as an orthogonal projection of $f\colon [t-\theta, t] \to \mathbb{R}$. Indeed, the trajectory of the input function $f$ could be reconstructed by the formula
% \begin{align*}
%     f(t- \theta') \approx \sum_{j=1}^N P_j\pr{\frac{2\theta'}{\theta}-1}c(t)_j.
% \end{align*}
% The HiPPO-LegS framework \citep{gu2020hippo} generalizes the LMU by projecting $f$ not onto a  fixed-length interval, but instead onto the whole domain $[0,t]$, conserving all the information of the trajectory of $f$. By solving \eqref{eq:intro_legs} yields
% \begin{align} \label{eq:LegS_projection}
%     c(t)_j &= \Big\langle f(s), P_{j-1}\pr{\frac{2s}{t}-1} \Big\rangle, \qquad 0 \leq s \leq t.
% \end{align}
% However, while LegS enjoys various empirical and theoretical advantages compared to LMU, 
% the $1/t$ factor present in \eqref{eq:intro_legs} introduces a singularity, which makes the analysis more involved. Compared to LTI systems such as the LMU, time-variant state-space models in machine learning is relatively underexplored, motivating the present study. 



\paragraph{Convergence analysis of SSMs from control theory.}
% {\color{red} What about these references? Maybe they belong in related works?}
State space models have been extensively studied in control theory \citep{kalman1960new, zabczyk2020mathematical}, with significant research dedicated to discretization schemes and their analysis \citep{kowalczuk1991discretization}. However, these results are not directly applicable to the LegS ODE due to their exclusive focus on LTI systems \citep{karampetakis2014error} or their assumption of discrete-time inputs \citep{meena2020discretization}. Furthermore, the objectives of state-space models in deep learning applications differ fundamentally from those in classical control theory, where for the latter, controlling or statistically estimating the state is usually the focus. This difference makes it challenging to directly adapt these results to modern deep-learning contexts.

% Other works from control theory dealing with singular or linear time-varying systems \citep{dewilde1998time, karampetakis2014error, meena2020discretization} are generally not applicable since the main objective of these works is to control the output, not to approximate the trajectory of the input function. 


% \subsection{Organization}
% % \textcolor{red}{Rewrite this after organizing the sections}

% Section \ref{sec:problem setting} introduces the problem setting and defines the notation that is used throughout the work. Section \ref{sec:well-posed} proves that the solution of the LegS ODE exists and is unique if the initial condition is fixed to a pre-determined value. Section \ref{sec:error_analysis} verifies the convergence of various discretization methods applied to the LegS ODE.  By reformulating the problem into a quadrature of $f(t)$, we prove convergence of all the discretization methods of interest for all Riemann integrable $f$.
% % Motivated by this view, we also suggest a speedup algorithm that could be used in all practical settings of HiPPO-LegS, with little or no degradation of performance. 

\section{Problem setting and preliminaries} \label{sec:problem setting}

% \paragraph{LegS ODE.}
In this work, we consider the LegS ODE
\begin{align}   \label{eq:legs} %\tag{LegS ODE}
    \derivative{c}(t) &= -\frac{1}{t}Ac(t) + \frac{1}{t}Bf(t) 
    % c_{k+1} &= (1-\frac{A}{k})c_k + \frac{1}{k}Bf_k\\
\end{align}
for $t\in (0,T]$, where $T>0$ is some terminal time and $c\colon [0,T]\rightarrow  \mathbb{R}^N$ is the state vector encoding the continuous-time history of the input function $f \colon  [0,T] \to \mathbb{R}$. The matrix $A \in \mathbb R^{N \times N}$ and vector $B \in \mathbb R^{N \times 1}$ are given by
\[
    \begin{aligned}
    A_{ij} &= \begin{cases}
        (2i-1)^{1/2}(2j-1)^{1/2} & \text{if} \quad i > j \\
        i & \text{if} \quad i=j \\
        0 & \text{if} \quad i < j ,
    \end{cases} \qquad\qquad
    B_j &= (2j-1)^{1/2}.
\end{aligned}
\]
Since $A$ is lower-triangular, we immediately recognize that $A$ is diagonalizable with simple eigenvalues $\{1, 2, \dots, N\}$. We denote the eigendecomposition as
\[
A=V DV^{-1},\qquad 
    D = \mathrm{diag}\,(1,2,\dots,N).
\]
with invertible $V\in\mathbb{R}^{N\times N}$.
In the indexing of $A_{ij}$, $B_j$, and $c_i(t)$, we have $i,j\in\{1,\dots,N\}$, i.e., we use $1$-based indexing. (The prior HiPPO paper \citep{gu2020hippo} uses $0$-based indexing.)

% {\color{red} Doesn't this $P$ notation conflict with the Legendre polynomials? Can we use $A=VD V^{-1} $ instead?}

% Throughout this paper, we choose the diagonalization scheme as $A = P^{-1}DP$, where $\pr{D}_{ij} = i\delta_{ij}$.





% We use $j \in J$ as the dimension-wise index of a vector, where $J = \{1, 2, \dots, N\}$. Specifically,
% \[
% \pr{c(t)}_j \quad \text{or} \quad c(t)_j
% \]
% refers to the $j$-th component of vector $c(t) \in \mathbb{R}^N$. If there is no confusion, it could also used to refer a vector that has the $\pr{c(t)}_j$ as the $j$-th component. Similarly,
% \[
% \pr{A}_{ij} \quad \text{or} \quad A_{ij}
% \]
% denotes the $i$-th row and $j$-th column of a matrix $A$. 


% \paragraph{Basic notation.}
% We use $e_n \in \mathbb{R}^N$ to refer to the $n$-th orthonormal standard basis, i.e., $\pr{e_n}_j = \delta_{nj}$, where $\delta$ is the dirac-delta notation.

\paragraph{Shifted Legendre polynomials.}
We write $P_j(x) \colon [-1,1] \to [-1,1]$ to denote the $j$-th Legendre polynomial, normalized such that $P_j(1)=1$, for $j=0,1,\dots$.
However, we wish to operate on the domain $[0,1]$, so we perform the change of variables $x \mapsto 2x-1$.
This yields, 
\[
\tilde{P}_j(x)=  P_j(2x-1) = \sum_{k=0}^j (-1)^j \binom{j}{k} \binom{j+k}{k} (-x)^k,
\]
the $j$-th \emph{shifted} Legendre polynomial, for $j=0,1,\dots$. The shifted Legendre polynomials satisfy the recurrence relation
\begin{align} \label{eq:shifted_legen_rel}
    x\tilde{P}_j'(x) = j\tilde{P}_j(x) + \sum_{k=0}^{j-1}(2k+1)\tilde{P}_k(x),
\end{align}
which can be derived by combining the following well-known identities \citep{arfken2011mathematical}
\begin{align*}
    (2n + 1)P_j(x) &= P_{j+1}'(x) - P_{j-1}'(x) \\
    P_{j+1}'(x) &= (n+1)P_j(x) + xP_j'(x).
\end{align*}


% We use $\dot{c}(t)$ or $\frac{d}{dt}c (t)$ interchangably to denote the derivative of functions. 



% For both well-posedness and discretization results, unless otherwise specified, we consider a compact domain $[0,T]$ on which the ODE is defined.
% For discrete-time analysis, we consider $n$ equidistant meshpoints starting from $t_0=0$, with timestep $h = \left\lfloor \frac{T}{n}\right\rfloor$. The exact solution at $t_n = nh$ is denoted as $c(t_n)$, while the numerical solution obtained by applying a discretization method at the $n$th step is denoted by $c_n$. 

% \subsection{LegS setting}
% We follow the LSI(Linear Scale Invariant) LegS formulation from the original HiPPO paper. In that case, the function is approximated by scaled Legendre Polynomials on a uniform measure $\mu_t = \left[ 0,t \right]$. Since the (length of the) interval is not fixed and vanishes at $t=0$, we investigate the mathematical soundness of it and look for feasible improvements.

% So the matrix $A$ and vector $B$ looks like: 
% \begin{align*}
%     A &= 
%     \begin{bmatrix}
%          1 & 0 & 0 \\
%          \sqrt{3} & 2 & 0  \\
%          \sqrt{5} & \sqrt{15} & 3\\
%          &&& \ddots
%     \end{bmatrix} \\
%     B &= \begin{bmatrix}
%         1 \\
%         \sqrt{3} \\
%         \sqrt{5} \\ 
%         \vdots
%     \end{bmatrix}
% \end{align*}

% The result could be derived by differentiating the measure and basis functions with respect to $t$, and recollecting the terms in vector form. 

% The original setting assumes that $f$ is a $C^1$-smooth, $L$-Lipschitz function $f:[0,\infty) \xrightarrow{} \mathbb R$. In this work, we only assume that $f$ is locally integrable and has a Lebesgue point at $t=0$ for the existence and uniqueness of solution proof. For error analysis, we assume $f$ is analytic.



% \label{subsec:disc_methods}

\paragraph{Numerical discretization methods.} 
In this work, we analyze the numerical methods of the LegS ODE used in the prior work \citep{gu2020hippo}. For all the discretization methodss, we consider a mesh grid with $n$ mesh points, with initial time $t_0=0$ and stepsize $h = T/n$. Starting from $c^0=c(0)$, we denote $k$-th step of the numerical method as $c^k$, and $f(kh)$ as $f^k$.

% Consider 
The \textbf{\underline{backward Euler}} method
\begin{align*}
    c^{k+1} &= \Big(I + \frac{1}{k+1}A\Big)^{-1}c^k + \Big(I + \frac{1}{k+1}A\Big)^{-1} \frac{1}{k+1}  B f^{k+1},
    \qquad\text{ for }k=0,1,2,\dots, n-1
\end{align*}
is well defined as is.
However, the \textbf{\underline{forward Euler}} method
\begin{align*}
    % c(t+\Delta t) - c(t) &\approx \Delta t [A_tc(t) + B_tf(t)]_{t=t} \\
    % c(t+\Delta t) &= (I + \Delta t A_t)c(t) + \Delta t B_tf(t) \\
    c^{k+1} &= \Big(I - \frac{1}{k}A\Big) c^k + \frac{1}{k}Bf^k,
    \qquad\text{ for }k=1,2,\dots, n-1
\end{align*}
and the \textbf{\underline{bilinear (trapezoidal)}} method
\begin{align*}
    c^{k+1} &= \Big(I + \frac{1}{2(k+1)}A\Big)^{-1} \Big(I - \frac{1}{2k}A\Big)c^k + \Big(I + \frac{1}{2(k+1)}A\Big)^{-1} \Big(\frac{1}{2k}f^k + \frac{1}{2(k+1)}f^{k+1}\Big)B,
    \qquad\text{ for }k=1,2,\dots, n-1
\end{align*}
are not well defined for $k=0$. One remedy would be to use the identity $c'(0) = \pr{A+I}^{-1} f'(0)$, which we derive in Lemma~\ref{lem:t=0 calc}. However, if $f$ is not differentiable at $t=0$, then even this remedy is not possible.
Hence, in Section~\ref{sec:error_analysis}, where we consider general $f$, we ``zero-out'' the ill-defined terms by setting them to be $0$.
So, for the \textbf{\underline{step $0$ of forward Euler}}, we set
\[
c^1 = c^0
    \qquad\text{ for }n=0,
\]
and for the \textbf{\underline{step $0$ of bilinear (trapezoidal)}}, we set
\begin{align*}
    c^{1} &= \pr{I + \frac{1}{2}A}^{-1} c^0 + \frac{1}{2} \pr{I + \frac{1}{2}A}^{-1}Bf^1,
    \qquad\text{ for }k=0.
\end{align*}



In the prior work \cite{gu2020hippo}, the authors sidestep the division by $0$ in the $1/k$ terms by shifting the $k$ index up by $1$, leading to the 
% \textbf{approximate forward Euler} method
% \begin{align*}
%     c_{n+1} &= \pr{I - \frac{1}{n+1}A}c_n + \frac{1}{n+1}Bf_{n+1} \qquad\text{ for }n=0,1,2,\dots
% \end{align*}
% and 
\textbf{\underline{approximate  bilinear}} method
\begin{align*}
    c^{k+1} &= \pr{I + \frac{1}{k+1}A/2}^{-1} \pr{I - \frac{1}{k+1}A/2}c^k + \pr{I + \frac{1}{k+1}A/2}^{-1} \pr{\frac{1}{k+1}f^{k+1}}B,\qquad\text{ for }k=0,1,2,\dots,n-1.
\end{align*}
In this work, we establish convergence of both the bilinear method (with zero-out) and the approximate bilinear method.

Lastly, prior work has also used the \textbf{\underline{Zero-order hold}} method
\begin{align*}
    c^{k+1} = e^{A\log\pr{\frac{k}{k+1}}}c^k + A^{-1} \pr{I - e^{A \log \pr{\frac{k}{k+1}}}}B f^k, \qquad\text{ for }k=0,1,2,\dots,n-1,
\end{align*}
where we set $e^{A\log\pr{\frac{n}{n+1}}}=0$ at $n=0$, consistent with the limit $e^{A\log\pr{\frac{n}{n+1}}}\rightarrow 0$ as $n\rightarrow 0^+$.
We also establish convergence for the zero-order hold method.


\paragraph{Convergence of numerical discretization methods.} 
% To numerically solve an initial value problem of the form 
% \begin{align*} 
%     c'(t) = g(t,c(t)) \label{eq:IVP} 
% \end{align*}
% for $t\ge 0$ with initial condition $c(0) = c^0$,
The numerical methods we consider are one-step methods of the form
\[
c^{k+1} = c^k + h\Phi(t_k, t_{k+1}, c^k, c^{k+1}; h), \qquad k=0,1,...,n-1
\]
with stepsize $h  =  T/n$ and $t_k=t_0+kh$ for $k=0,\dots,n-1$, approximating the solution to the initial value problem
\begin{align*} 
    c'(t) = g(t,c(t)). 
\end{align*}
Here, $\Phi$ is a numerical integrator making the approximation
\[
\Phi(t_k, t_{k+1}, c^k, c^{k+1}; h)\approx c(t_{k+1})-c(t_k)= \int^{t_k+1}_{t_k}g(s,c(s))\;ds.
\]

To analyze such methods, one often estimates the \emph{local truncation error} (LTE) $T_k$ at timestep $t_k$ as
\begin{align*}
T_k 
&= \frac{c(t_{k+1}) - c(t_k)}{h} - \Phi(t_k, t_{k+1}, c(t_k), c(t_{k+1}); h)
\end{align*}
and then estimates its accumulation to bound the \emph{global error}
\[
e_n = c(t_n) - c^n
\]
which is calculated at the endpoint.
We say that a numerical discretization method is \emph{convergent} if the global error converges to $0$, i.e., if
\[
\|c(t_n) - c^n\|\rightarrow 0,\qquad\text{ as }\,n\rightarrow\infty.
\]
Further, we quantify the \emph{convergence rate} with the \emph{order} of the method: we say the method has order $p$ if 
\begin{align*}
    \| c^n - c(t_n) \| \leq \mathcal{O}\pr{1/n^p},\qquad\text{ as }\,n\rightarrow\infty.
\end{align*}
Classical ODE theory states that if the right-hand-side $g$ in the initial value problem is continuous with respect to $c$ and $t$, and Lipschitz continuous with respect to $c$, the solution exists and is unique in an interval including the initial point $t_0=0$. Moreover, under the same conditions, the global error can be bounded with the local truncation error \citep{ascher_petzold_1998,Süli_Mayers_2003}. However, this standard theory does not apply to the LegS ODE due to the singularity at $t=0$, and the non-smoothness of the input function $f$.


% To obtain a numerical solution, we consider a meshgrid with $n$ meshpoints defined on the interval $[0,T]$. 
% % and a discretization method, one would want to obtain an estimate of the numerical error introduced by the approximation. At each mesh point $t_n = t_0 + nh$, one performs numerical integration
% In this case, we denote the input function at step $t_n$ as $f(t_n) = f(nh) = f_n$. 




% (one-step method) where $h = \left\lfloor \frac{T}{n} \right\rfloor $ is the stepsize and $\Phi$ is the numerical integrator of interest. Naturally, one would expect the discretization error to decrease as the stepsize gets smaller (with more computational cost) and converge to 0 as $h \to 0$. 


% \textbf{Global error}: \[
% e_n = y(t_n) - y_n
% \]
% \textbf{Truncation error}: \begin{align*}
% T_n 
% % &= \frac{1}{h} \Big( (y(t_{n+1})-y(t_{n})) - (y_{n+1}-y_n) \Big)
% &= \frac{y(t_{n+1}) - y(t_n)}{h} - \Phi(t_n, y(t_n); h)
% \end{align*}



\paragraph{Absolute continuity on a half-open interval.} 
For $T\in (0,\infty)$, we say a function $c\colon(0,T]\rightarrow\mathbb{R}^N$ is absolutely continuous if its restriction to the closed interval $[\varepsilon,T]$ for any $\varepsilon\in (0,T)$ is absolutely continuous. (Recall that the standard definition of absolute continuity assumes a closed interval for the domain.) Even if $\lim_{t\rightarrow 0^+}c(t)$ is well defined and finite, the continuous extension of $c$ to $[0,T]$ may not be absolutely continuous on $[0,T]$. In other words, absolute continuity on $[\varepsilon,T]$ for all $\varepsilon\in (0,T)$ does not imply absolute continuity on $[0,T]$. We discuss this technicality further in Section~\ref{sec:well-posed}, in the discussion following Theorem~\ref{thm:exist and unique}.

\paragraph{Lebesgue point.}
Let $f\colon [0,T]\rightarrow\mathbb{R}$ be Lebesgue measurable and integrable. 
We say $f$ has a Lebesgue point at $t=0$ if
\[
\lim_{\varepsilon\rightarrow 0^+}\frac{1}{\varepsilon}\int^\varepsilon_0|f(s)-f(0)|\;ds=0.
\]
If $f(t)$ is continuous at $t=0$, then $f$ has a Lebesgue point at $t=0$.

\section{LegS is well-posed} \label{sec:well-posed}
In this section, we show that the LegS ODE is well-posed despite the singularity. Crucially, however, we show that there is no freedom in choosing the initial condition. 

% \js{Should we provide the definition of a.c. on half interval here?}

% \subsection{Existence and uniqueness of solution}

% The following theorem states that the LegS equation is well-posed.

% \begin{theorem}[Existence and uniqueness] \label{thm:exist and unique}
% {\color{red} Shouldn't $c$ be absolutely continuous? The given solution is not necessarily differentiable.}
%     For $M>0$, we say $c$ is a solution if $c \in C([0,M);\mathbb R^N)$ is differentiable on $(0,M)$ and satisfies \eqref{eq:legs} for $t \in (0,M)$ and $c(0) = c_0$.
%     % For $M>0$, we say $c$ is a solution if $cC^0([0,M) is differentiable ;\mathbb R)$, satisfies \eqref{eq:1d_LegS} for $t \in (0,M)$ and $c(0) = c_0$.
%     % Consider the conditions
%     % \begin{itemize}
%     %     \item [(i)] For all $t>0$, $\int_{0}^{t} f(s) ds$ is well-defined.
%     %     \item [(ii)] Limit $\lim_{t\to0} f(t)$ exists.
%     %     \item [(iii)] $c_0 = \lim_{t\to0} f(t)$.
%     % \end{itemize}
%     % If (i), (ii), (iii) are satisfied, the solution exists and is unique. Furthermore, if (iii) is not satisfied, the solution does not exist.
%     Assume $f$ is Lebesgue measurable and locally integrable. 
%     Further, assume $f$ is continuous at $t=0$.
%     % Assume locally integrable $f(t)$ 
%     % has a Lebesgue point at $t=0$. 
%     % is continuous at $t=0$.
%     Then, the solution exists and is unique if $c_0 = A^{-1}Bf(0)$. Otherwise, the solution does not exist.
%     {\color{red} Why not include this here? $A^{-1}B = e_1 = [1, 0, ..., 0]^{t}$}
% \end{theorem}

\begin{theorem}[Existence and uniqueness] \label{thm:exist and unique}
% {\color{red} Shouldn't $c$ be absolutely continuous? The given solution is not necessarily differentiable.}
    For $T>0$ and $c_0 \in \mathbb{R}^N$, we say $c\colon [0,T]\rightarrow \mathbb R^N$ is a solution (in the extended sense) of the LegS ODE if $c$ is
    continuous on $[0,T]$,
    absolutely continuous on $(0,T]$, $c$ satisfies \eqref{eq:legs} for almost all $t \in (0,T]$, and $c(0) = c_0$.
    % For $M>0$, we say $c$ is a solution if $cC^0([0,M) is differentiable ;\mathbb R)$, satisfies \eqref{eq:1d_LegS} for $t \in (0,M)$ and $c(0) = c_0$.
    Assume $f\colon [0,T]\rightarrow \mathbb R$ is Lebesgue measurable, integrable, and 
    % Further, assume $f$ is continuous at $t=0$.
    % Assume locally integrable $f(t)$ 
    has a Lebesgue point at $t=0$. 
    % is continuous at $t=0$.
    Then, the solution exists and is unique if $c_0 = f(0)e_1$, where $e_1\in \mathbb{R}^N$ is the first standard basis vector.
    Otherwise, if $c_0 \ne  f(0)e_1$, a solution does not exist.
\end{theorem}

\begin{proof}
    Since $A$ is diagonalizable, the problem can be effectively reduced to $N$ $1$-dimensional problems. 
    % Using the usual spectral decomposition, let $A = VDV^{-1}$ with $D$ being a diagonal matrix. 
    Recall $A = VDV^{-1}$ where $D=\mathrm{diag}\,(1,2,\dots,N)$. 
    We see the LegS ODE \eqref{eq:legs} could be rewritten as 
    \begin{align*}
        \derivative{c}(t) 
        &= -\frac{1}{t} A c(t) + \frac{1}{t} B f(t) %\\&
        = -\frac{1}{t} VDV^{-1} c(t) + \frac{1}{t} B f(t). 
    \end{align*}
    Multiply both sides by $V^{-1}$ and denote $\tilde{c}(t) = V^{-1}c(t)$. Then, the ODE becomes
    \begin{align*}
        \derivative{\tilde{c}}(t) = V^{-1}\derivative{c}(t) &= -\frac{1}{t} D V^{-1} c(t) + \frac{1}{t} V^{-1}B f(t) = -\frac{1}{t} D\tilde{c}(t) + \frac{1}{t} V^{-1}B f(t),
    \end{align*}
    % and hence
    % \begin{align*}
    %     \derivative{\tilde{c}}(t)&= -\frac{1}{t} D\tilde{c}(t) + \frac{1}{t} V^{-1}B f(t)
    % \end{align*}
    which is a decoupled ODE with respect to $\tilde{c}$. Recalling $D_{jj} = j$, we see that the $j$-th component of the above equation is
    % Considering the $j$-th component of the above equation yields
    \begin{align}   
        \derivative{\tilde{c}}_j(t)&= -\frac{j}{t}\tilde{c}_j(t) + \frac{d_j}{t} f(t) \label{eq:1d_legs}
    \end{align}
    where $d_j = (V^{-1}B)_j$ and $\tilde{c}_j$ is $j$-th component function of $\tilde{c}$. % and noting that $(D)_{jj} = j$. 
    Since $V$ and $V^{-1}$ are absolutely continuous bijective maps from $\mathbb{R}^N$ to $\mathbb{R}^N$, 
    % the existence and uniqueness of $c(t) \in \mathbb{R}^N$ is satisfied if and only if $\tilde{c}_j(t) \in \mathbb{R}$ exists and is unique for all %$j = 1, 2, ..., N$. 
    the existence and uniqueness of the solution of the LegS ODE is satisfied if and only the existence and uniqueness of the solution of the ODE \eqref{eq:1d_legs} is satisfied for all 
    $j \in \{1, 2, ..., N\}$. 

    We now proceed by examining the existence and uniqueness of the solution 
    % $\tilde{c}_j \in C([0,T];\mathbb{R})$ 
    of the ODE \eqref{eq:1d_legs}. 
    % that is absolutely continuous on $(0,T]$, satisfying \eqref{eq:1d_legs} for almost all $t \in (0, T]$. 
    We first establish existence by presenting the explicit form of the solution.  
    Define $\tilde{c}_j\colon[0,T]\to\mathbb{R}$ as
    \begin{equation}    \label{eq:closed_form_solution}
        \tilde{c}_j(t) = \begin{cases}
            \frac{d_j}{t^j} \int^t_0 s^{j-1} f(s) ds  & \text{if} \quad t \in (0,T] \\
            \frac{d_j}{j} f(0)  & \text{if} \quad t = 0.
        \end{cases}
    \end{equation}
    By fundamental theorem of calculus, $\frac{d}{dt} \pr{ \int^t_0 s^{j-1} f(s) ds } = t^{j-1} f(t)$ holds for almost all $t\in(0,T]$ and thus $\tilde{c}_j$ is differentiable for almost all $t\in(0,T]$. 
    Therefore,
    \[
        t^j \pr{ \derivative{\tilde{c}}_j(t) + \frac{j}{t} \tilde{c}_j(t) } 
        % = t^j \derivative{\tilde{c}}_j(t) + j t^{j-1} \tilde{c}_j(t) 
        = \frac{d}{dt} (t^j\tilde{c}_j(t)) 
        = \frac{d}{dt} \pr{ d_j \int^t_0 s^{j-1} f(s) ds }
        = d_j t^{j-1} f(t) 
    \]
    holds for almost all $t\in(0,T]$. 
    Dividing both sides by $t^j$, we conclude $\tilde{c}_j$ satisfies \eqref{eq:1d_legs} for almost all $t\in(0,T]$.
    
    We now show $\tilde{c}_j$ is continuous on $[0,T]$. 
    It is sufficient to check $\tilde{c}_j$ is continuous at $t=0$ by showing
    % Since $f$ is continuous at $t=0$, from L'H\^{o}pital's rule we have
    % \begin{equation*}
    %     \begin{aligned}
    %         \lim_{t\to0^+} \tilde{c}_j(t)
    %         &= \lim_{t\to0^+} \frac{d_j}{t^j} \int^t_0 s^{j-1} f(s)ds  \\ 
    %         &= \lim_{t\to0^+} \frac{d_j }{j t^{j-1}} \pr{ t^{j-1} f(t) }
    %         = \frac{d_j}{j} \lim_{t\to0^+} f(t) 
    %         = \frac{d_j}{j} f(0) = \tilde{c}_j(0).
    %     \end{aligned}
    % \end{equation*}
    % The last equation comes from the continuity of $f$ at $t=0$. 
    \begin{align*}
        \lim_{t \to 0} \frac{d_j}{t^j} \int^t_0 s^{j-1} f(s) ds &= \frac{d_j}{j}f(0).
    \end{align*}
    % So basically we want to show that assuming $\tilde{c}_j(t)$ is 
    % \begin{equation*}
    %         \tilde{c}_j(t) = 
    %             \frac{d_j}{t^j} \int^t_0 s^{j-1} f(s) ds  \quad \text{if} \quad t \in (0,T]
    % \end{equation*}
    % that $\tilde{c}_j(0)= \frac{d_j}{j}f(0)$. Drop $d_j$ for now. 
    Since $f$ is locally integrable and has a Lebesgue point at $t=0$, observe that
    \begin{equation*}
    % \label{eq:fact_from_Lebesgue_point}
        0 = \lim_{t \to 0} 
        \frac{1}{t} \int_0^t |f(s)-f(0)|ds  
        = 
        \lim_{t \to 0} 
        \int_0^1 |f(tx) - f(0)|dx,
    \end{equation*}
    where the second equality follows by change of variables $x = s/t$. Therefore, we could deduce
    % For $t>0$, observe
    % \begin{equation*}
    %     \abs{ \tilde{c}_j(t) - \frac{1}{j} f(0) }
    %     = \abs{ \frac{1}{t} \int^t_0 \frac{1}{t^{j-1}} s^{j-1} f(s) ds - \frac{1}{t} \int_{0}^{t} \frac{1}{j} f(0) ds }
    %     \le \frac{1}{t} \int_{0}^{t} \abs{ \frac{1}{t^{j-1}} s^{j-1} f(s) - \frac{1}{j} f(0) } ds.
    % \end{equation*}
    \begin{align*}
        \abs{\lim_{t \to 0} \frac{d_j}{t^j} \int^t_0 s^{j-1} f(s) ds - \frac{d_j}{j}f(0)    
        } &\leq
        \lim_{t \to 0} 
        \frac{d_j}{t} \int_0^t \abs{ \frac{1}{t^{j-1}} s^{j-1} f(s) - \frac{1}{j} f(0)} ds \\
        &= \lim_{t \to 0} 
        d_j \int_0^1 \abs{x^{j-1} f(tx) - \frac{1}{j}f(0) } dx \\
        &\leq \lim_{t \to 0} 
        d_j \int_0^1 x^{j-1} |f(tx) - f(0)| dx  + d_j \int_0^1 \abs{x^{j-1} - \frac{1}{j}} |f(0)| dx \\
        &\leq \lim_{t \to 0} 
        d_j \int_0^1 |f(tx) - f(0)| dx  
        + d_j |f(0)| \int_0^1 \abs{x^{j-1} - \frac{1}{j} }  dx \\
        &= 0,
    \end{align*}
    % where the third inequality comes from the fact $x\le1$ and the followed equality comes from \eqref{eq:fact_from_Lebesgue_point}. 
    concluding that $\tilde{c}_j$ is continuous at $t=0$. 
    Lastly, $\tilde{c}_j$ is absolutely continuous on $(0,T]$, since for every $[t_0,t] \subset (0,T]$, 
    % $\derivative{\tilde{c}}_j$ exists almost everywhere and $\tilde{c}_j(t) = \tilde{c}_j(t_0) + \int_{t_0}^{t} \derivative{\tilde{c}}_j(s) ds$ holds.  
    both $\frac{1}{t^j}$ and $\int^t_0 s^{j-1} f(s) ds$ are absolutely continuous on $[t_0,t]$ and therefore their product is also absolutely continuous on $[t_0,t]$. 
    % \textcolor{red}{XXDRAFTSTART absolute cont at $t=0$}
    % $f$ is not absolutely continuous even if it is continuous at $t=0$. Counterexample exists. 
    % \textcolor{red}{XXDRAFTEND absolute cont}
    Hence we conclude that $\tilde{c}_j$ is a solution of the ODE \eqref{eq:1d_legs}. 

    We now establish uniqueness. 
    Suppose $\hat{c}_j$ is another solution of the ODE \eqref{eq:1d_legs}. 
    Multiplying both sides of \eqref{eq:1d_legs} by %$\exp(j\log{t}) = t^j$ 
    $t^j$ and reorganizing, for almost all $t\in(0,T]$ we have
    % \begin{align*}
    %     \frac{d}{dt} (t^jc(t)) 
    %     &= d_j t^{j-1} f(t) \\
    %     t^j c(t) - t_0^j c(t_0) &= d_j \int_{t_0}^t s^{j-1} f(s) ds
    % \end{align*}
    \[
        \frac{d}{dt} (t^j\hat{c}_j(t)) 
        = t^j \pr{ \derivative{\hat{c}}_j(t) + \frac{j}{t} \hat{c}_j(t) } 
        = d_j t^{j-1} f(t) .
    \]
    Since $\hat{c}_j$ is a solution, it is absolutely continuous on $(0,T]$, therefore $t^j\hat{c}_j(t)$ is absolutely continuous on $(0,T]$. 
    Thus for $[t_0,t] \subset (0,T]$, 
    % integrating from $t_0$ to $t$ we obtain
    by fundamental theorem of calculus we obtain
    \[
        t^j \hat{c}_j(t) - t_0^j \hat{c}_j(t_0) = d_j \int_{t_0}^t s^{j-1} f(s) ds.
    \]
    Taking limit $t_0 \to 0^+$, since $\hat{c}_j$ is a solution it is continuous at $0$, we have
    \[
        t^j \hat{c}_j(t) = d_j \int_{0}^t s^{j-1} f(s) ds.
    \]
    Dividing both sides by $t^j$ we conclude
    \begin{equation*}
       \hat{c}_j(t) = \frac{d_j}{t^j} \int^t_0 s^{j-1} f(s)ds = \tilde{c}_j(t) %\label{eq:test func}
    \end{equation*}
    for all $t \in (0,T]$. 
    % Furthermore, from L'H\^{o}pital's rule we have
    It remains to check $\hat{c}_j(0)=\tilde{c}_j(0)$. 
    Since $\hat{c}$ is continuous at $0$, we know $\hat{c}_j(0)=\lim_{t\to0^+}\hat{c}_j(t)$. Thus
     \begin{equation*}
        \hat{c}_j(0)
        = \lim_{t\to0^+} \hat{c}_j(t)
        = \lim_{t\to0^+} \tilde{c}_j(t)
        = \tilde{c}_j(0). 
    \end{equation*}
    % Again from L'H\^{o}pital's rule, we have
    % \begin{equation*}
    %     \lim_{t\to0^+} \hat{c}_j(t)
    %     % = \lim_{t\to0^+} \frac{d_j}{t^j} \int^t_0 s^{j-1} f(s)ds 
    %     % = \lim_{t\to0^+} \frac{d_j }{j t^{j-1}} \pr{ t^{j-1} f(t) }
    %     = \frac{d_j}{j} f(0) 
    % \end{equation*}
    % The last equation comes from the continuity of $f$ at $t=0$. 
    Therefore, we conclude $\hat{c}_j(t) = \tilde{c}_j(t)$ for all $t\in[0,T]$, 
     the solution of the ODE \eqref{eq:1d_legs} is unique. 
    
    % Now we continuously extend this function to $t=0$. Since $f$ is continuous at $t=0$, we can immediately apply L'H\^{o}pital's rule. 

    % Otherwise, 
    % \textcolor{blue}{Try not to overuse arrows (for limits)}
    % \begin{align*}
    %     \left| \hat{c}(t) - \frac{d_j}{j} f(0) \right|
    %     &= \left|\frac{d_j}{t^j} \int_0^t s^{j-1} f(s)ds - \frac{d_j}{j} f(0) \right| \\
    %     &= d_j \left| \frac{1}{t^j} \int_0^{t^j} \frac{1}{j}f(u^{1/j})du - \frac{1}{j}f(0) \right|
    %     % & \leq \frac{d_j}{j} \frac{1}{t^j} \int_0^{t^j} |f(u^{1/j}) - f(0)| du \\
    %     \xrightarrow[t\to 0]{}0
    % \end{align*}
    % since $t=0$ is a Lebesgue point of $f$ (\textcolor{red}{Insufficient argument}).
    
    % Hence we let $\hat{c}(0)_j = \frac{d_j}{j}f(0)$, and $\hat{c}(t)_j$ is a solution by construction. 
    % Observe that $\lim_{t\to0}\tilde{c}(t) = d_j f(0)$ since $\left| \frac{d_j}{t}\int_0^tf(s)ds - f(0) \right| \le \frac{d_j}{t}\int_0^t \left| f(s)- f(0) \right| ds \xrightarrow[t\to0]{}0$. Then by construction, this $\tilde{c}(t)$ is a solution.

    % To organize, 
    % \begin{equation*}
    %     \derivative{\tilde{c}}(t) 
    %     =  -\frac{1}{t} D\tilde{c}(t) + \frac{1}{t} V^{-1}B f(t)
    % \end{equation*}
    
    % We proceed by showing uniqueness. Assume $c_1, c_2 \in C^0[0,M)$ differentiable on $(0,M)$ satisfy (\ref{eq:1d_legs}) for $t \in (0, M)$. Then, 
    % \begin{align*}
    %     \frac{d}{dt} (c_1 - c_2)(t) &= -\frac{j}{t}(c_1-c_2)(t)\\
    %     c_1(t)-c_2(t) &= \frac{C}{t^{j}}
    % \end{align*}
    % where $C$ is a constant. However, continuity at $t=0$ enforces $C = 0$. Hence the only feasible solution is $\hat{c}(t)_j$ the one we have calculated. Moreover, this unique function has a fixed candidate for its value at $t=0$.

    As a result, we conclude the solution of the LegS ODE uniquely exists if $\tilde{c}_j(0) = \frac{d_j}{j}f(0)$, and it is given by $c = V \tilde{c}$. 
    Finally, we show the unique solution $c = V \tilde{c}$ should satisfy $c(0) = f(0)e_1$.  
    % Folding $\tilde{c}_j(0)$ back to the vector form, 
    From %from $V^{-1}c(0)_j = \tilde{c}_j(t) = \frac{d_j}{j}f(0) = \frac{1}{j}(V^{-1}B)_j f(0)$,
    \[
        V^{-1}c_j(0) = \tilde{c}_j(0) = \frac{d_j}{j}f(0) = \frac{1}{j}(V^{-1}B)_j f(0),
    \]
    % so if $c_0 =  d_j f(0)$, the solution is unique, and otherwise there is no solution.
    we see
    \[
        V^{-1}c(0)  = \text{diag} \pr{ 1, \frac{1}{2}, \frac{1}{3}, ... , \frac{1}{N} } (V^{-1}B) f(0) = D^{-1} V^{-1} B f(0).
    \]
    Multiplying both sides by $V$, we conclude
    \begin{align*}
        c(0) 
        = V D^{-1} V^{-1} B f(0)
        = ( V^{-1} D V )^{-1}Bf(0)
        % {\color{red} ??? XXX}
        = A^{-1}B f(0) 
        = f(0)e_1
    \end{align*}
    % \begin{align*}
    %     V^{-1}c(0) 
    %     &=  \text{diag}(1, 1/2, 1/3, ... ) (V^{-1}B) f(0) \\
    %     c(0) &= \underbrace{P^{-1} \text{diag}(1, 1/2, 1/3, ... ) P}_{= A^{-1}} B f(0)\\
    %     &= A^{-1}B f(0) \\
    %     &= f(0)e_1
    % \end{align*}
    % Note that this is obtained through an efficient calculation ($\mathcal{O} (N^2)$), since $A$ is a lower-triangular matrix. 
    where $e_1 = [1, 0, \dots, 0] ^{t}$.  
    Therefore if $c_0 = f(0)e_1$, the solution exists and is unique, and otherwise, there is no solution.
\end{proof}



\begin{remark}    
Theorem~\ref{thm:exist and unique} does not guarantee that $c$ is absolutely continuous on the closed interval $[0,T]$, only on the half-open interval $(0,T]$.
    The following lemma provides a counterexample of a continuous input function $f$ such that the corresponding solution of the LegS ODE is not absolutely continuous on \( [0,T] \).    
\end{remark}

\begin{lemma}
Let $T=1/2$.
Consider the LegS ODE with $f\colon[0,T]\to\mathbb{R}$ defined as
\begin{equation*}
    f(t) =
    \begin{cases}
         \frac{d}{dt} \pr{ \frac{t^2}{\log(1/t)} \sin \pr{ 1/t }}  & \text{ if } 0< t \le1/2 \\
        0 & \text{ otherwise}.
    \end{cases}
\end{equation*}
Since $f$ is continuous on $[0,T]$ (including at $t=0$) it satisfies the conditions of Theorem~\ref{thm:exist and unique}.
However, the solution of the LegS ODE is not absolutely continuous on \( [0,T] \).
% However, the solution of LegS ODE is not absolutely continuous on $[0,\delta]$ for every $\delta>0$.

\end{lemma}


\begin{proof}
    We first show $f$ satisfies the desired conditions. 
    Since $f$ is continuous on $(0,1/2]$, it is suffices to show $f$ is continuous at $t=0$. 
    Calculating the differentiation, we see that
    \begin{equation*}
        f(t)= 
        \frac{d}{dt}  \pr{ \frac{t^2}{\log(1/t)} \sin(1/t) }
        % = \frac{t \left(2 \log(1/t)+1\right) }{\log^2(1/t)} \sin(1/t)
        % - \frac{1}{\log(1/t)} \cos(1/t) 
        = \frac{2t}{\log(1/t)} \sin(1/t)
        + \frac{t}{\log^2(1/t)} \sin(1/t)
        - \frac{1}{\log(1/t)} \cos(1/t) 
    \end{equation*}
    holds for $t\in(0,1/2]$. 
    Since $\lim_{t\to0^+} \log(1/t) = \lim_{u\to\infty} \log(u) = \infty$, it follows that $\lim_{t\to0^+} f(t) = 0$. 
    % \begin{equation*}
    %     \lim_{t\to0+} f(t)
    %     =  \lim_{t\to0+} \frac{t}{\log(1/t)} \left(2 + \frac{1}{\log(1/t)} \right) \sin(1/t)
    %     - \lim_{t\to0+} \frac{1}{\log(1/t)} \cos(1/t)
    %     = 0.
    % \end{equation*}
    Therefore, $f$ is continuous at $t=0$, so it is integrable and has a Lebesgue point at $0$. 

    Now, we show that $c(t)$ is not absolutely continuous on $[0,T]$. It is sufficient to prove that the first component $c_1(t)$ is not absolutely continuous on $[0,T]$. 
    % From the closed form of the solution provided in \eqref{eq:closed_form_solution}, we know 
    % \begin{equation*}
    %     c_1(t) = \begin{cases}
    %         \frac{d_1}{t} \int_{0}^{t} f(x) dx & \text{ if } t>0 \\
    %         0 & \text{ if } t=0.
    %     \end{cases}
    % \end{equation*}
    Plugging the definition of $f(t)$ to \eqref{eq:closed_form_solution}, we obtain
    \begin{equation*}
        c_1(t) = \begin{cases}
            c_1\pr{ \frac{1}{2} } & \text{ if } t>1/2 \\
            d_1 \frac{t}{\log(1/t)} \sin(1/t) & \text{ if } t\in (0,1/2] \\
            0 & \text{ if } t=0.
        \end{cases}
    \end{equation*}
    Let $\delta>0$ be arbitrary positive number. 
    For $N>0$ such that $\frac{1}{(2N+1) \pi + \frac{\pi}{2}} < \min\set{\delta,\frac{1}{2},T}$,  consider
    \begin{equation*}
        t_n = \begin{cases}
           \pr{ (2N+n+\frac{1}{2})\pi }^{-1} & \text{ if } n \text{ is odd } \\
            \pr{ (2N+n)\pi }^{-1} & \text{ if } n \text{ is even} .
        \end{cases}
    \end{equation*}
    Define $x_n = t_n$ and $y_n=t_{n+1}$ for $n\ge1$. 
    Then $\sum_{n=1}^{\infty} \pr{ y_{n} - x_{n} } = t_1
         = \frac{1}{(2N+1) \pi + \frac{\pi}{2}}< \delta$.
    % \begin{equation*}
    %     \sum_{n=1}^{\infty} \pr{ y_{n} - x_{n} }
    %     = 
    %      \sum_{n=1}^{\infty} \pr{ t_{n} - t_{n+1} }
    %      = t_1
    %      = \frac{1}{(2N+1) \pi + \frac{\pi}{2}}< \delta.
    % \end{equation*}
    However, since
    \begin{equation*}
        \abs{ c_1(x_n) - c_1(y_{n}) } 
        = \abs{ c_1(t_n) - c_1(t_{n+1}) }
        = \begin{cases}
            \abs{c_1(t_n)} & \text{ if $n$ is odd }   \\
            \abs{c_1(t_{n+1})} & \text{ if $n$ is even}   ,
        \end{cases}
    \end{equation*}
    we have
    \begin{equation*}
        \begin{aligned}
            \sum_{n=1}^{\infty} \abs{ c_1(x_n) - c_1(y_{n}) } 
            % &= \sum_{n=1}^{\infty} \abs{ c_1(t_n) - c_1(t_{n+1}) }  \\
            &= \abs{c_1(t_1)} + 2\sum_{m=1}^{\infty} \abs{ c_1(t_{2m+1})} \\
            &> 2 d_1 \sum_{m=1}^{\infty} \frac{t_{2m+1}}{\log(1/t_{2m+1})} %\\&
            = 2 d_1 \sum_{m=1}^{\infty} \frac{1}{ 
(2N + 2m + 3/2)\pi \log((2N + 2m + 3/2)\pi)} \\
            &> 2 d_1 \sum_{m=1}^{\infty} \frac{1}{ 
2\pi(N + m + 1) \log(2\pi(N + m + 1))} 
            = 2 d_1 \sum_{m=N+2}^{\infty} \frac{1}{ 
            2m\pi \log(2m\pi)} 
            = \infty.
        \end{aligned}
    \end{equation*}
    % since $\sum_{n=2}^{\infty} \frac{1}{n \log n} = \infty$. 
    The last equality follow from integral test:
    \begin{equation*}
        \sum_{m=N+2}^{\infty} \frac{1}{ 
            2m\pi \log(2m\pi)}
        \ge \int_{N+2}^{\infty} \frac{1}{2\pi t \log (2\pi t)} dt
        = \frac{1}{2\pi} \int_{2\pi(N+2)}^{\infty} \frac{1}{u \log u} du 
        = [\log(\log(u))]_{2\pi(N+2)}^{\infty}
        = \infty. 
    \end{equation*}
    Since $(x_n,y_n) \subset [0,T]$ are disjoint and $\delta>0$ was arbitrary, we conclude $c_1$ is not absolutely continuous on $[0,T]$.
\end{proof}


% \textcolor{red}{$f$ is not absolutely continuous on $[0,t]$ given the setup in Theorem~\ref{thm:exist and unique}. Provide a counterexample (for continuous $f$) here?}

% \textcolor{red}{Moreover, we can prove that if $f$ is differentiable at $t=0$, then $c$ is absolutely continuous on $[0,T]$. However, it is not sure if this is a tight condition. Should we add this discussion?}

\begin{remark}
   Recall that the motivation of the LegS ODE is to provide an online approximation of the input function $f$. 
   By change of variables, one can show that $\{ \frac{\sqrt{2j-1}}{t} P_{j-1}\pr{\frac{2s}{t} -1} \}_{j \in \mathbb{N}}$ is an orthogonal basis on the interval $[0, t]$, with respect to the $L^2([0,t])$ norm. 
    The following corollary shows that the solution found in Theorem \ref{thm:exist and unique} is the projection of $f$ onto this basis.
\end{remark}


\begin{corollary} \label{cor:sol_legen_form}
    The solution $c$ of the LegS ODE as defined as in Theorem \ref{thm:exist and unique} is, if it exists, an $L^2$-approximation of $f$ on $\frac{1}{t}\mathbb{I}_{[0,t]}$ for all $t \in [0,T]$ in the sense that the $j$-th component of $c(t) \in \mathbb{R}^N$ is given by
    \begin{align} \label{eq:sol_legen_form}
        c_j(t) &= \frac{1}{t} \Big\langle f(\cdot ), \sqrt{2j-1} P_{j-1}\big(\tfrac{2\,\cdot} {t}-1\big) \Big\rangle_{L^2([0,t])}
    =
    \frac{\sqrt{2j-1}}{t} \int_0^t f(s) P_{j-1}\big(\tfrac{2s} {t}-1\big) \;ds,
    \end{align}
    for all $t \in (0,T]$, where $P_{j-1}$ denotes the $(j-1)$-th Legendre polynomial. 
\end{corollary}

\begin{proof}
    If we assume that the solution exists, we know by the previous theorem that $c(0) = A^{-1}B f(0)$. 
    Further, in the proof of Theorem \ref{thm:exist and unique}, we have that $c(t) = V\tilde{c}(t)$ for all $t \in (0,T]$ where
    \begin{align*}
        \tilde{c}_j(t) 
        &= \frac{\pr{V^{-1}B}_j}{t^j} \int_0^t s^{j-1}f(s)ds.
    \end{align*}
    The $j$-th component of $c$ could be expressed as 
    \begin{align}
        c_j(t) &= \sum_{k=1}^N V_{jk}
        \frac{\pr{V^{-1}B}_k}{t^k} \int_0^t s^{k-1}f(s)ds  %\nonumber \\&
        = \frac{1}{t} \int_0^t 
        \underbrace{{\sum_{k=1}^N V_{jk} \pr{V^{-1}B}_k x^{k-1}}}_{=G_j(x)}
        f(s)ds \label{eq:diag_sol}
    \end{align}
    where we made the substitution $x = \frac{s}{t}$. Since other terms do not depend on the index $j$, we simplify this expression by analyzing $G_j(x)$, regarding as $x$ as a symbolic (differentiable) variable. Expression $G_j(x)$ could be rewritten as 
    \begin{align*}
        G_j(x) &= \sum_{k=1}^N V_{jk} \pr{V^{-1}B}_k x^{k-1} \\
        &= \sum_{k,l=1}^N V_{jk} x^{k-1} V^{-1}_{kl} B_l \\
        &= e_j^t V \text{diag} \pr{1, x, \dots, x^{j-1}, \dots, x^{N-1}} V^{-1}B.
    \end{align*}
    Define $G(x)$ as a length $N$ vector with its $j$-th component being $G_j(x)$, which is a polynomial of $x$.
    Then we obtain a matrix differential equation with respect to the symbolic variable $x$ :
    \begin{align*}
        x \frac{dG}{dx}(x) = \pr{A-I} G(x).
    \end{align*}
    Since we know the precise structure of matrix $A$, letting $g_{j-1}(x) = \frac{G_j(x)}{\sqrt{2j-1}} $, we can obtain the following recurrence relation for $g_j$'s, 
    \begin{align*}
        xg_{j-1}'(x) = (j-1)g_{j-1} + \sum_{k=0}^{j-2}(2k+1)g_k(x),
    \end{align*}
    which is precisely the recurrence relation \eqref{eq:shifted_legen_rel} for the shifted Legendre polynomials. Since $G_j(1) = B_j = \sqrt{2j-1}$, we conclude that $g_j$ is equal to the $(j-1)$-th shifted Legendre polynomial $\tilde{P}_{j-1}$. 
    Finally, incorporating this observation into \eqref{eq:diag_sol}, we can rewrite the solution of the LegS ODE as 
    \begin{align*}
        c_j(t) &= \frac{1}{t}\int_0^t \sqrt{2j-1} \tilde{P}_{j-1}\pr{x} f(s)ds %\\&
        = \frac{\sqrt{2j-1}}{t} \int_0^t P_{j-1} \pr{\frac{2s}{t}-1} f(s)ds
    \end{align*}
    for all $t \in (0, T]$.
\end{proof}


% With stronger guarantees of $f$ at $t=0$, we can obtain the one-sided differentiability of $c(t)$ at $t=0$, as shown in the following Lemma. 

\begin{remark}
    The well-posedness argument of Theorem~\ref{thm:exist and unique} crucially relies on the fact that all eigenvalues of $A$ are positive.
    To see what happens when $A$ has negative eigenvalues, consider the case $N=1$ and $A=-1$. 
    This leads to the ODE
    \begin{align*}
        \frac{d}{dt}c(t) &= \frac{1}{t} c(t) + \frac{1}{t} f(t),\qquad c(0)=c_0,
    \end{align*}
    which has solutions
    \begin{align*}
        c(t) = t \int_0^t \frac{1}{s^2}f(s)ds + Ct
    \end{align*}
    for any $C\in \mathbb{R}$. Since the initial condition does not determine the value of $C$, the solution is not unique.
\end{remark}


% \paragraph{Ill-posedness with negative eigenvalues.}
\begin{remark}
    If a stronger condition, such as the (one-sided) differentiability of $f(t)$ at $t=0$ is provided, the derivative of $c(t)$ at $t=0$ can be calculated as in the following lemma. In setups where $f'$ is available, this identity could be used to implement the first iteration of the forward Euler method and the bilinear method.
    % This result provides a means to correctly implement the first iteration of the forward Euler method and the bilinear method. However, it is worth noting that in many practical settings using the HiPPO formulation to deep state space models, the derivative of the input function $f$ is not available. 
\end{remark}

\begin{lemma} [Behavior at $t=0$] \label{lem:t=0 calc}
% {\color{red}
% Consider the setup of XXX.}
Consider the setup of Theorem~\ref{thm:exist and unique}, and further assume that $f$ is one-sided differentiable at $t=0$.
Then, the one-sided derivative $c'(0)$ exists and is given by
        \begin{align}
        c'(0) 
        % &= \underbrace{P^{-1} \text{diag}(1/2, 1/3, 1/4, ... ) P}_{= (A+I)^{-1}} B f'(0)\\
        &= (A+I)^{-1}B f'(0).\label{eq:c'(0)}
    \end{align}
\end{lemma}
\begin{proof}    
We examine the differentiability of $V^{-1}c(t) = \tilde{c}(t) = (\tilde{c}_j(t))_{j=1}^N$ by checking for each component. Recall that for the $j$-th component we have $\tilde{c}_j(t) = \frac{d_j}{t^j}\int_0^t s^{j-1}f(s)ds$ and $\tilde{c}_j(0) = \frac{d_j}{j}f(0)$. 
We will denote the one-sided derivative of $c(t)$ and $f(t)$ at $t=0$ as $c'(0)$ and $f'(0)$, respectively. 
Then,
\begin{align*}
    \lim_{t \to 0^+} \frac{\tilde{c}
_j(t) - \tilde{c}_j(0)}{t} 
    &= \lim_{t \to 0^+} d_j \frac{\frac{1}{t^j} \int_0^t s^{j-1} f(s)ds - \frac{1}{j}f(0)}{t} \\
    &= \lim_{t \to 0^+} d_j \frac{\int_0^t s^{j-1} f(s)ds - \frac{t^j}{j}f(0)}{t^{j+1}} \\
    &\overset{(1)}{=} \lim_{t \to 0^+} d_j \frac{t^{j-1}f(t) - t^{j-1} f(0)}{(j+1)t^{j}} %\\&
    = \lim_{t \to 0^+} \frac{d_j}{j+1} \frac{f(t) - f(0)}{t} %\\&
    = \frac{d_j}{j+1} f'(0)
\end{align*}
where L'H\^{o}pital's rule was used for $(1)$. 
 Folding back to vector form, recalling $V^{-1}c'_j(0) = \tilde{c}'_j(t) = \frac{d_j}{j+1}f'(0) = \frac{1}{j+1}(V^{-1}B)_j f'(0)$, we obtain
    \begin{align*}
        V^{-1}c'(0) 
        &=  \text{diag} \pr{ \frac{1}{2}, \frac{1}{3}, ... , \frac{1}{N+1} } %\text{diag}(1/2, 1/3, 1/4, ... ) 
        (V^{-1}B) f'(0) 
    \end{align*}
and hence $c'(0) = \pr{A+I}^{-1}B$. 
Note that $(A+I)^{-1}B = [1/2, 1/(2\sqrt{3}), 0, ..., 0]^{t}$.
\end{proof}



% \subsection{Well-posedness relies on the structure of $A$}

% This kind of ODE would be relevant if $A$ had any negative eigenvalues. In this case, we find that the solution is not unique. One could easily verify that the following function 

% solves the ODE for any $C \in \mathbb{R}$. However, regardless of the initialization $c(t_0)=c_0$ given at $t_0=0$, the solution will always have the ambiguity of the $Ct$ term. 



% This fact is actually necessary for the well-posedness of the ODE. Here we explain why.

% For illustrative purposes, consider the $N=1$ case. Then, the solution for the ODE 
% \begin{align*}
%     \dot{c}(t) &= -\frac{1}{t} c(t) + \frac{1}{t} f(t)
% \end{align*}
% is
% \begin{align*}
%     c(t) &= \frac{1}{t} \int_a^t f(s)ds
% \end{align*}
% where $a$ is decided by the initial condition. However, the inclusion of $t=0$ in the domain forces $a=0$. (or $\int_0^a f(s)ds = 0$)



\section{Convergence of LegS discretization schemes} \label{sec:error_analysis}

% \textcolor{blue}{Section composition.}
% For the remainder of the paper, we shift gears and aim to answer the following question:

% \begin{center}
%     \textit{When applying the discretization methods introduced in Section \ref{sec:problem setting} to the LegS ODE, does the numerical solution $c_n$ converge to the exact solution $c(t)$ as the stepsize $h$ converges to 0?}    
% \end{center}

In this section, we address the convergence of the numerical discretization methods introduced in Section~\ref{sec:problem setting}, i.e., do the methods produce numerical solutions $c^n$ that converge to the exact continuous-time solution $c(t)$ as $h\rightarrow 0$?

The answer is yes, but, as we discuss in Section~\ref{subsec:smooth_f}, the standard analysis based on local truncation error does not lead to a convergence guarantee for all of the schemes under consideration, and such approaches would require certain local regularity conditions on $f$, such as (Lipschitz) continuity. Rather, in Section~\ref{subsec:legs is quadrature}, we identify the numerical schemes as quadrature rules on the input function $f$. Using this insight, in Section~\ref{subsec:riemann_f}, we show that the discretization schemes are convergent under the general assumption of Riemann integrability of $f$.


% The extension to allow general Riemann integrable $f$ is important due to the nature of the application. The HiPPO memory unit is utilized in deep learning to analyze sequence data, such as audio sequence, and there is no a prior reason to expect such input signals to be smooth. In fact, many signals of interest are expected to have discontinuities, so it is meaningful to have a mathematical guarantee that the method is sound for such data.

Extending the framework to accommodate general Riemann integrable functions $f$ is important, given the nature of the application. The HiPPO memory unit is used in deep learning to analyze sequence data, such as audio signals. For such data, there is no inherent expectation of smoothness, and discontinuities are to be expected. Therefore, we aim to guarantee that the mathematics remains sound for such data.






% This extension is important XXX
% Note that this is an important point since in practice, the input function is not assumed to be smooth, and in many cases it is not even guaranteed to be continuous. For example, when used for sequential models such as RNNs, the input function $f$ is defined as the hidden state of the model, which is defined on discrete-time. The ideal way to model these types of systems would be to use step functions, which are not continuous.


% It is a well-known fact that if the ODE is non-singular in the sense that it satisfies the conditions in Picard's existence and uniqueness theorem, the global error converges with the same rate as the truncation error. 
% While 
% the LegS ODE \eqref{eq:legs} does not satisfy that criterion due to the singularity at $t=0$, 
% it could be proven that the global error converges to $0$ as the stepsize decreases if all the truncation errors uniformly converge to $0$. 


\subsection{Convergence for smooth $f$} \label{subsec:smooth_f}
Discretization methods of ODEs with well-behaved right-hand-sides have a well-established theory based on the local truncation error (LTE), and the standard techniques can be applied to the LegS ODE despite the singularity at $t=0$. For example, it can be shown that LTE for the forward Euler method applied to the LegS ODE satisfies
\[
|T_k| \leq \frac{1}{2} h M_2, \qquad M_2 = \max_{t \in [0, T]} |c''(t)|.
\]
However, when $f$ is not differentiable, then $c''(t)$ may not be bounded, and this approach, as is, fails to yield a convergence guarantee.

% The proof follows from a relatively straightforward extension of the standard arguments.


% Hence, to prove convergence using this bound, it should be verified that $c''(t)$ is well-defined and bounded.
% If $f$ is sufficiently smooth, i.e. twicely continuously differentiable, it is straightforward to show that $c''$ is bounded on a compact domain. 
% However, if $f$ is not continuously differentiable, we cannot assume that $c''(t)$ is bounded. This could be illustrated by the simple dimension $N=1$ case; i.e., the exact solution for the 1D LegS equation is
% \begin{align*}
%     c(t) = \frac{1}{t}\int_0^t f(s)ds.
% \end{align*}
% It could be easily shown that if $f(t)$ is not differentiable, $c''(t)$ could be ill-defined. 

Another issue for the approximated bilinear method is that the LTE does not converge to $0$. 
For $N=1$, the approximated bilinear method reduces to 
\begin{align*}
    c^{k+1}
    &= \frac{2k+1}{2k+3}c^k + \frac{2}{2k+3}f^{k+1} = c^k - \frac{2}{2k+3}(c^k - f^{k+1}).
\end{align*}
% Hence, we can solve this recurrence relation by plugging $c^0 = f_0$ to obtain $c^n = \frac{1}{2n+1}f_n + \frac{2}{2n+1}\sum_{k=1}^n f_k$. 
Using the exact solution $c(t) = \frac{1}{t}\int_0^t f(s)ds$, the exact value of the LTE is
\begin{align*}
    hT_k
    &= c(t_{k+1}) - c(t_{k}) + \frac{2}{2k+3} (c(t_k) - f^{k+1}) \\
    &= \frac{1}{(k+1)h}\int_0^{(k+1)h}f(s)ds - \frac{2k+1}{k(2k+3)h}\int_0^{kh}f(s)ds - \frac{2}{2k+3}f^{k+1} .
    % &= \frac{1}{h^2} \int^h_0 f(s)ds - \frac{1}{3h}f(0) -\frac{2}{3h}f(h)
    % &= \frac{1}{h} \int^h_0 f(s)ds - \frac{1}{2}( f(0) + f(h)) -\frac{1}{6}( f(h) - f(0))
\end{align*}
With the linear function $f(x) = ax$ as a particular choice, we obtain that at step $n=0$,
% \begin{align*}
%     T_0 &= \frac{a}{2}(n+1) - \frac{n(2n+1)}{2(2n+3)}a - \frac{2a}{2n+3}(n+1) = -\frac{a(n+1)}{2(2n+3)}\rightarrow -\frac{a}{4}\ne 0.
% \end{align*}
\begin{align*}
    T_0 &= \frac{a}{2} - \frac{2a}{3} 
    = -\frac{a}{6} \ne 0.
\end{align*}
Thus, the LTE of the approximated bilinear method does not vanish as $h\rightarrow 0$.
Consequently, a naive global error analysis based on the LTE will not guarantee convergence. In Section~\ref{subsec:riemann_f}, we employ an alternative proof technique to establish convergence.



% This is an important gap since it implies that the use of the approximated bilinear method in the original work lacks sufficient theoretical justification. Fortunately, in the next section we will prove that, although not with the same order of accuracy, the approximated bilinear method converges to the exact solution as $h \to 0$. 

% Second, the proof fails also if we do not have sufficient smoothness $f$. This is because the truncation error is bounded by Taylor's theorem, which requires high-order derivatives of $c(t)$. Since $c(t)$ is a (weighted) integral of $f(t)$, $c^{(n)}(t)$ may not be bounded or well-defined on $[0,t]$ without the smoothness of $f(t)$. 



% \section{LegS is a quadrature of $f$} \label{sec: legs is quadrature}


\subsection{LegS discretizations are quadratures of $f$} \label{subsec:legs is quadrature}
In this section, we provide the key insight that we can identify the discretization methods as quadrature rules on the input function $f$. 
Recall that the LegS ODE was proposed for online approximation of the input function $f(t)$ on the interval $[0,t]$.
Specifically, Corollary \ref{cor:sol_legen_form} says that the $j$-th component of the solution is given by
\begin{align*}
c_j(t) &= \frac{\sqrt{2j-1}}{t} \int_0^t P_{j-1} \Big(\frac{2s}{t} - 1 \Big) f(s) ds,
\end{align*}
where $P_{j-1}$ denotes the $(j-1)$-th Legendre polynomial, for $j=1,2,\dots,N$.
% {\color{red}[ XXX provide citationXXX ]}.
So $c_j(t)$ is a (signed) weighted integral of $f(\cdot)$ on $[0,t]$. The following lemma shows that the numerical schemes of Section~\ref{sec:problem setting} can, in fact, be interpreted as quadrature rules with uniformly spaced nodes, and in Section~\ref{subsec:riemann_f}, we show that these quadratures approximate the integral.


% We can verify this is a solution for \eqref{eq:legs} by direct computation, and it is the unique solution by Theorem \ref{thm:exist and unique}.
% Here, the obvious but important observation is that since $c(t)$ is a weighted average of $f$ on the interval $[0,t]$, numerically estimating $c(t)$ is equivalent to estimating the integral of $f(t)$. 
% \textbf{In other words, the discretization of the LegS ODE(with respect to $c(t)$) is equivalent to the (numerical) quadrature of $f(t)$.} We proceed to formalize this idea by presenting the following lemma. 

% A point of subtlety worth noting is that, as discussed in Section \ref{subsec:disc_methods} and Section \ref{subsec:t=0 calc}, $\dot{c}(0)$ is not well-defined if $f$ is not differentiable at $t=0$. Since we are considering general input functions $f$, we simply skip undefined iterations, e.g. $c_1 = c_0$ for forward Euler methods. As will be shown, this modification does not affect the convergence results. 

\begin{lemma} \label{lem:numsol_linear}
Consider applying any of the discretization  methods introduced in Section~\ref{sec:problem setting}
    (forward Euler, backward Euler, bilinear, approximate bilinear, or zero-order hold)
    to the LegS ODE \eqref{eq:legs}, with initial time $t_0=0$ and timestep $h = T/n$. 
    % Then, the numerical solution $c^n$ at step $n$ satisfies
    % \begin{align*}
    %     c^n &= \frac{1}{n} \sum_{l=0}^n \alpha^{(n)}_l f^l= \frac{1}{n} \sum_{l=0}^n \alpha^{(n)}_l f(lh),
    % \end{align*}
    % for all $n = 1,2, \dots, n$, where $\alpha^{(n)}_l \in \mathbb{R}^N$ only depends on $l$ and $n$. 
    Then, the numerical solution $c^n$ at step $n$ can be expressed in the form
    \begin{align*}
        c^n &= \frac{1}{n} \sum_{l=0}^n \alpha^{(n)}_l f^l = \frac{1}{n} \sum_{l=0}^n \alpha^{(n)}_l f(lh),
    \end{align*}
    for all $n = 1, 2, \dots, n$, for some $\alpha^{(n)}_l \in \mathbb{R}^N$ that depend only on $l$ and $n$.
\end{lemma}


\begin{proof}
    % We first prove for the forward Euler method. 
    % We defer the proof for other methods to Appendix \ref{sec:appendixB}.

    % For notational simplicity, we define the following quantities for $n \in \mathbb{N}$, $n \geq 2$ for $Q_n$:
    For notational simplicity, we define $Q_n$ for $n \geq 2$, and $\tilde{Q}_n$, $R_n$, $\tilde{R}_n$ for $n \geq 1$ as follows:
    % \textcolor{red}{Should change notation $Q_n \to Q^n$?}
    % \begin{align*}
    %     Q_n &:= \prod_{j=1}^{n-1} \pr{I - \frac{1}{j+1}A} \\
    %     \tilde{Q}_n &:= \prod_{j=1}^n \pr{I - \frac{1}{2j}A}\\
    %     R_n &:= \prod_{j=1}^n \pr{I + \frac{1}{j}A}^{-1} \\
    %     \tilde{R}_n &:= \prod_{j=1}^n \pr{I + \frac{1}{2j}A}^{-1}.
    % \end{align*}
    \begin{align*}
        Q_n &:= \prod_{j=1}^{n-1} \pr{I - \frac{1}{j+1}A}, \qquad
        \tilde{Q}_n := \prod_{j=1}^n \pr{I - \frac{1}{2j}A}, \\
        R_n &:= \prod_{j=1}^n \pr{I + \frac{1}{j}A}^{-1}, \qquad\,\,\,\,
        \tilde{R}_n := \prod_{j=1}^n \pr{I + \frac{1}{2j}A}^{-1}.
    \end{align*}
    We use the $\prod$ notation when the multiplications are commutative. Note that all $Q_n, \tilde{Q}_n, R_n, \tilde{R}_n$ are invertible. 
    
    We start by proving for the forward Euler method. 
    Recall that the forward Euler method yields the following recurrence relation at step $n$:
\begin{align*}
    c^{n+1} &= \pr{ I - \frac{1}{n}A}c^n + \frac{1}{n} B f^n.
\end{align*}
Repeating this procedure, we can obtain an exact formula for the numerical solution obtained by applying forward Euler method to the LegS ODE. By induction we obtain,
\begin{align*}
    c^{n+1} 
    &= \pr{ I - \frac{1}{n}A }c^n + \frac{1}{n}Bf^n \\
    &= \pr{ I - \frac{1}{n}A} \pr{ I - \frac{1}{n-1}A} c^{n-1} + \pr{ I - \frac{1}{n}A} \frac{1}{n-1}Bf^{n-1} + \frac{1}{n}Bf^n \\
    &= Q_n \pr{c^1 + Bf^1} + Q_n \sum_{l=2}^n \frac{1}{l} Q_{l}^{-1} B f^l.
\end{align*}
As explained in Section~\ref{sec:problem setting}, we `zero out' the ill-defined iteration, thereby letting $c^1=c^0$. Hence we have
\begin{align} \label{eq:forward_numsol}
    c^{n}  
    = \frac{1}{n} \sum_{l=0}^n \alpha^{(n)}_l f(lh)
    = Q_{n-1} (e_1 f^0 + B f^1 ) +  Q_{n-1} \sum_{l=2}^{n-1} \frac{1}{l} Q_l^{-1} B f^l.
\end{align}

One can verify that $\alpha_l^{(n)}$ depends only on $l$ and $n$. 

For the backward Euler method, we start from
\begin{equation*}
\begin{aligned}
    c^{n+1} 
    &= \pr{ I + \frac{1}{n+1} A  }^{-1} c^{n}
     + \pr{ I + \frac{1}{n+1} A  }^{-1} \frac{1}{n+1} B f^{n+1}.
\end{aligned}
\end{equation*}

Then, we can derive inductively 
\begin{equation*}
    \begin{aligned}
    c^{n+1} 
    &= \pr{ I + \frac{1}{n+1}A }^{-1} c^n + \pr{ I + \frac{1}{n+1}A }^{-1} \frac{1}{n+1}Bf^{n+1} \\
    &= \pr{ I + \frac{1}{n+1}A }^{-1} \pr{ I + \frac{1}{n}A }^{-1} c^{n-1} + \pr{ I + \frac{1}{n+1}A }^{-1} \pr{ I + \frac{1}{n}A }^{-1} \frac{1}{n}Bf^n + \pr{ I + \frac{1}{n+1}A }^{-1} \frac{1}{n+1}Bf^{n+1} \\
    % &\cdots \\
    &= R_{n+1} (c^0 + Bf^1) + R_{n+1} \sum_{l=1}^{n} \frac{1}{l+1} R_{l}^{-1} B f^{l+1}.
    \end{aligned}
\end{equation*}
Thus we can verify $c^n$ has the desired form with the specific expression
\begin{equation} \label{eq:back_numsol}
    \begin{aligned}
    c^n 
    = \frac{1}{n} \sum_{l=0}^n \alpha^{(n)}_l f(lh)
    =    R_{n} (e_1f^0 + Bf^1) + R_{n} \sum_{l=1}^{n-1} \frac{1}{l+1} R_{l}^{-1} B f^{l+1}.
    \end{aligned}
\end{equation}

For the bilinear method, we start from
\begin{align*}
    c^{n+1} &= \pr{I + \frac{1}{n+1}A/2}^{-1}\pr{I - \frac{1}{n}A/2}c^n +\pr{I + \frac{1}{n+1}A/2}^{-1} \frac{1}{2} \pr{\frac{1}{n}f^n + \frac{1}{n+1}f^{n+1}}B.
\end{align*}

Similarly, by induction, we obtain
\begin{equation*}
    \begin{aligned}
    c^{n+1} 
    &= \pr{ I + \frac{1}{n+1}A/2 }^{-1} \pr{ I - \frac{1}{n}A/2 } c^n + \pr{ I + \frac{1}{n+1}A/2 }^{-1} \frac{1}{2} \pr{\frac{1}{n}Bf^{n} + \frac{1}{n+1}Bf^{n+1} }\\
    &= \pr{ I + \frac{1}{n+1}A/2 }^{-1} \pr{ I + \frac{1}{n}A/2 }^{-1} \pr{ I - \frac{1}{n}A/2 } \pr{ I - \frac{1}{n - 1}A/2 } c^{n-1}\\
    &\phantom{=}+ \pr{ I + \frac{1}{n+1}A/2 }^{-1} \pr{ I + \frac{1}{n}A/2 }^{-1}\pr{ I - \frac{1}{n - 1}A/2 } \frac{1}{2} \pr{\frac{1}{n-1}Bf^{n-1} + \frac{1}{n}Bf^{n} }\\
    &\phantom{=}+ \pr{ I + \frac{1}{n+1}A/2 }^{-1} \frac{1}{2} \pr{\frac{1}{n}Bf^{n} + \frac{1}{n+1}Bf^{n+1} }\\
    &= \tilde{Q}_n \tilde{R}_{n+1} \pr{I+A/2} c^1 + 
    \tilde{Q}_{n} \tilde{R}_{n+1} \sum_{l=1}^{n} \tilde{Q}_l^{-1} \tilde{R}_l^{-1}  \frac{1}{2}\pr{\frac{1}{l}Bf^l + \frac{1}{l+1}Bf^{l+1}}.
    \end{aligned}
\end{equation*}

As for the forward Euler case, we `zero out' the ill-defined term in the first iteration. This yields $c^1 = \pr{I+\frac{A}{2}}^{-1}c^0 + \pr{I+\frac{A}{2}}^{-1}\pr{\frac{1}{2}f^1}$. 
Then,
\begin{equation*} 
    c^n =
    % \tilde{Q}_{n-1}\tilde{R}_{n}e_1f^0
    % +\tilde{Q}_{n-1}\tilde{R}_{n}\frac{f^1}{2}
    \tilde{Q}_{n-1}\tilde{R}_{n} \pr{e_1f^0
    +\frac{f^1}{2}}
    +\tilde{Q}_{n-1}\tilde{R}_{n} \sum_{l=1}^{n-1} \tilde{Q}_l^{-1} \tilde{R}_l^{-1}  \frac{1}{2}\pr{\frac{1}{l}Bf^l + \frac{1}{l+1}Bf^{l+1}}.
\end{equation*}
Rearranging terms, 
\begin{align}\label{eq:bilin_numsol}
    c^n 
    = \frac{1}{n} \sum_{l=0}^n \alpha^{(n)}_l f(lh)
    &=
    \tilde{Q}_{n-1}\tilde{R}_{n}e_1f^0
    + \tilde{Q}_{n-1}\tilde{R}_{n}
    \sum_{l=1}^{n-1} \frac{1}{2l}
    \pr{
    \tilde{Q}_l^{-1} \tilde{R}_l^{-1} 
    + \tilde{Q}_{l-1}^{-1} \tilde{R}_{l-1}^{-1} 
    }Bf^l
    + \tilde{R}_n \tilde{R}_{n-1}^{-1} \frac{1}{2n} B f^n
\end{align}
where we define $\tilde{Q}_0=\tilde{R}_0=I$. Thus we recover the desired form for $c^n$. 

For the approximate bilinear method, we start from
\begin{align*}
    c^{n+1} &= \pr{I + \frac{1}{n+1}A/2}^{-1}\pr{I - \frac{1}{n+1}A/2}c^n + \pr{I + \frac{1}{n+1}A/2}^{-1} \pr{\frac{1}{n+1}f^{n+1}}B.
\end{align*}

By induction, we obtain
\begin{equation*}
    \begin{aligned}
    c^{n+1} 
    &= \pr{ I + \frac{1}{n+1}A/2 }^{-1} \pr{ I - \frac{1}{n+1} A/2 } c^n + \pr{ I + \frac{1}{n+1}A/2 }^{-1} \pr{\frac{1}{n+1}Bf^{n+1} }\\
    &= \pr{ I + \frac{1}{n+1}A/2 }^{-1} \pr{ I + \frac{1}{n}A/2 }^{-1} \pr{ I - \frac{1}{n+1}A/2 } \pr{ I - \frac{1}{n}A/2 } c^{n-1}\\
    &\phantom{=}+ \pr{ I + \frac{1}{n+1}A/2 }^{-1} \pr{ I + \frac{1}{n}A/2 }^{-1}\pr{ I - \frac{1}{n+1}A/2 } \pr{\frac{1}{n}Bf^{n} }\\
    &\phantom{=}+ \pr{ I + \frac{1}{n+1}A/2 }^{-1} \pr{\frac{1}{n+1}Bf^{n+1} }\\
    &= \tilde{Q}_{n+1} \tilde{R}_{n+1} c^0 + 
    \tilde{Q}_{n+1} \tilde{R}_{n+1} \sum_{l=1}^{n} \tilde{Q}_{l+1}^{-1} \tilde{R}_l^{-1} \pr{\frac{1}{l+1}Bf^{l+1}}.
    \end{aligned}
\end{equation*}

Thus we can verify $c^n$ has the desired form with the specific expression
\begin{equation} \label{eq:apbil_numsol}
    \begin{aligned}
    c^n &=
    \tilde{Q}_{n}\tilde{R}_{n}c^0 + \tilde{Q}_{n}\tilde{R}_{n} \sum_{l=1}^{n-1} \tilde{Q}_{l+1}^{-1} \tilde{R}_l^{-1} \pr{\frac{1}{l+1}Bf^{l+1}}.
    \end{aligned}
\end{equation}

For the Zero-order hold method, we start from
\begin{align*}
    c^{n+1} = e^{A\log\pr{\frac{n}{n+1}}}c^n + \pr{I - e^{A \log \pr{\frac{n}{n+1}}}} A^{-1} B f^n.
\end{align*}
By induction, we obtain
\begin{align*}
    c^{n+1} &= e^{A\log\pr{\frac{n}{n+1}}} e^{A\log\pr{\frac{n-1}{n}}}c^{n-1} 
    + \pr{I-e^{A\log\pr{\frac{n}{n+1}}}} A^{-1}Bf^n
    + e^{A\log\pr{\frac{n}{n+1}}} \pr{I - e^{A\log\pr{\frac{n-1}{n}}}} A^{-1}Bf^{n-1} \\
    &= e^{A\log\pr{\frac{n-1}{n+1}}}c^{n-1}
    + \pr{I-e^{A\log\pr{\frac{n}{n+1}}}} A^{-1}Bf^n
    +  \pr{e^{A\log\pr{\frac{n}{n+1}}} - e^{A\log\pr{\frac{n-1}{n+1}}}} A^{-1}Bf^{n-1} \\
    &= \sum_{k=0}^{n} \pr{e^{A\log\pr{\frac{n+1-k}{n+1}}} - e^{A\log\pr{\frac{n-k}{n+1}}}}A^{-1}Bf^{n-k}.
\end{align*}
Note that in this expression, we are denoting (with abuse of notation) $e^{\log{0}}=0$.
Thus we can verify $c^n$ has the desired form with the specific expression
\begin{align} \label{eq:zoh_numsol}
    c^n = \frac{1}{n} \sum_{l=0}^n \alpha^{(n)}_l f(lh) = \sum_{l=0}^{n-1} \pr{e^{A\log\pr{\frac{l+1}{n}}} - e^{A\log\pr{\frac{l}{n}}}} A^{-1}B f^l.
\end{align}
\end{proof}

% The result of Lemma \ref{lem:numsol_linear} implies that the discretized numerical solution for the LegS ODE is expressed as a linear combination of $f_l = f(lh)$'s. 
% This is due to the fact that the expression for $\dot{c}(t)$ is $f$-linear and $c$-linear, which is a property of state-space representations. This does not require the system to be LTI, allowing us to leverage it for our analysis.


\subsection{Convergence for Riemann integrable $f$} \label{subsec:riemann_f}

In this section, we prove the convergence of all discretization methods of interest for Riemann integrable $f$'s. 
% This is the most general condition we could impose on $f$, since we cannot expect for any quadrature method to converge for all Lebesgue integrable functions. 
% {\color{red} wasn't the bilinear method also considered in the previous result?}
In particular, we prove the convergence of the approximate bilinear method, justifying its use for the experiments in the HiPPO paper \citep{gu2020hippo}. 
The results are summarized in the following theorem:


\begin{theorem} [Convergence of discretization schemes for Riemann integrable $f$]
\label{thm:conv_riemann}
    Consider the LegS equation \eqref{eq:legs} with dimension $N\ge 1$ and domain $t \in [0,T]$, where $T>0$. 
    Assume $f$ is Riemann integrable on $[0,T]$.
    Let $n \ge 1$ be the number of mesh points and let $h=T/n$.
    Consider discretization methods with initialization
    \[
    % c^0_j = \delta_{1j} f(0) e_1
    c^0 = f(0) e_1
    \]
    % for $j=1,\dots, N$,
    % where $\delta_{1j}$ is the Kronecker delta, 
    and iterations
    \begin{align*}
    c^{k+1} &= c^k + h \Phi(t_k,t_{k+1},c^k,c^{k+1}; h) 
    \end{align*}
    using mesh points $t_k = kh$ for $k = 0, 1, 2, \dots, n$, where $\Phi$ is the one-step integrator defined by the given discretization method. Denote the exact solution at step $n$ as $c(t_n) = c(nh) \in \mathbb{R}^N$. Then, for all the forward Euler, backward Euler, bilinear, approximate bilinear, and zero-order hold methods defined in Section \ref{sec:problem setting}, we have convergence of the numerical solution to the exact solution in the sense that
\[
% \sup_{n=0,1,\dots, n}
\|c^n-c(t_n)\|\rightarrow 0,\qquad\text{ as } \quad n \rightarrow \infty.
\]
\end{theorem}

In light of Lemma \ref{lem:numsol_linear}, we can denote the iterates of the numerical schemes as $c^n = \frac{1}{n} \sum_{l=0}^n \alpha^{(n)}_l f^l$. If we can show
\[
c^n_j = \frac{1}{n} \sum_{l=0}^n \pr{\alpha^{(n)}_l}_jf^l 
\rightarrow
c_j(t) = \frac{\sqrt{2j-1}}{t}\int_0^t P_{j-1}\pr{\frac{2s}{t}-1}f(s)ds,
\]
as $n \to \infty$, and where $P_j$ is the $j$-th Legendre polynomial, we are done.

% We briefly summarize our proof strategy here. First, we observe that, for each component $j$ of the solution, if $\pr{\alpha^{(n)}_l}_j \to P_{j-1}\pr{\frac{2l}{n}-1}$ as $n \to \infty$ for all $l \in \{1, 2, \dots, n-1\}$, the numerical solution converges to the exact solution. Next, for fixed $n$, we consider an interpolating function $F^{(n)}_j$ that interpolates the $n-1$ points: $\pr{\frac{1}{n}, \pr{\alpha^{(n)}_1}_j}, \pr{\frac{2}{n}, \pr{\alpha^{(n)}_2}_j}, \dots, \pr{1-\frac{1}{n}, \pr{\alpha^{(n)}_{n-1}}_j}$. 
% Lastly, we will conclude the proof by showing that this function sequence $\{F^{(n)}_j\}_{n \in \mathbb{N}}$ converges to a limit $\sqrt{2j-1}\tilde{P}_{j-1}$.

Instead of directly characterizing the coefficients $\alpha^{(n)}_l$, the key idea is to consider a function sequence defined on $[0,1]$ that interpolates those points. This significantly reduces the complexity of analyzing the asymptotic behavior of the numerical solution as the number of mesh points $n$ goes to infinity. 
We start with
an elementary lemma that enables this approach.

\begin{lemma} \label{lem:converge_condition}
    Let $f \colon [0,t] \to \mathbb{R}$ be a Riemann integrable function.
    Let $\{G^{(n)}\}_{n \in \mathbb{N}}$ be a sequence of continuous functions defined on $[0,1]$ uniformly converging to $G \in C[0,1]$.
    Then, for $h = t/n$,
    \begin{align*}
        \frac{1}{n} \sum_{l=1}^n G^{(n)} \pr{\frac{l}{n}} f\pr{lh}
        \to
        \frac{1}{t}\int_0^t G \pr{\frac{s}{t}}f(s)ds
    \end{align*}
    as $n \to \infty$. 
\end{lemma}

\begin{proof}
    Fix $\epsilon > 0$. Since $f$ is Riemann integrable and $G$ is continuous, $G\pr{\frac{s}{t}}f(s)$ is Riemann integrable for $s \in [0,t]$. Hence we can find $N_1 \in \mathbb{N}$ such that for all $n \geq N_1$, $\left| \frac{1}{n} \sum_{l=1}^n G \pr{\frac{l}{n}}f\pr{lh} - \frac{1}{t}\int_0^t G \pr{\frac{s}{t}}f\pr{s} ds \right| < \epsilon / 2$. Since $f$ is bounded, let let $\sup_{x \in [0,t]} |f(x)| \leq M$. Then we can find $N_2 \in \mathbb{N}$ such that for all $n \geq N_2$, $\|G - G^{(n)}\| < \frac{\epsilon}{2M}$, due to the uniform convergence of $G^{(n)}$. Then, by choosing $n \in \mathbb{N}$ such that $n \geq \max\{N_1,N_2\}$ :
    % Since $f$ is Riemann integrable and $G$ is continuous, $G\pr{\frac{s}{t}}f(s)$ is Riemann integrable for $s \in [0,t]$. Then, the Riemann sum $\frac{1}{n} \sum_{l=1}^n G \pr{\frac{l}{n}}f\pr{lh}$ converges to $\frac{1}{t}\int_0^t G \pr{\frac{s}{t}}f\pr{s} ds$ as $n \to \infty$. Since $f$ is bounded, let $\sup_{x \in [0,t]} |f(x)| \leq M$, and
    \begin{align*}
        \left| \frac{1}{t}\int_0^t G \pr{\frac{s}{t} }f(s)ds - \frac{1}{n}\sum_{l=1}^n G^{(n)} \pr{\frac{l}{n}} f \pr{lh} \right| 
        &\leq \left| \frac{1}{t}\int_0^t G \pr{\frac{s}{t}} f(s)ds - \frac{1}{n} \sum_{l=1}^n G \pr{\frac{l}{n}}f\pr{lh} \right| \\
        &+ \frac{1}{n} \sum_{l=1}^n \left| G \pr{\frac{l}{n}}f\pr{lh} - G^{(n)} \pr{\frac{l}{n}} f\pr{lh} \right| \\
        &\leq \frac{\epsilon}{2} + \sup_{x \in [0,t]} \left| f(x) \right| \max_{l \in \{0, \dots, n\}} \left| G\pr{\frac{l}{n}} - G^{(n)} \pr{\frac{l}{n}} \right| \\
        &\leq \frac{\epsilon}{2} + M \| G - G^{(n)} \|_{\sup} \leq \epsilon
    \end{align*}
\end{proof}

This result is relevant since if we interpret the shifted Legendre polynomials as $G\pr{\frac{s}{t}}$ in the lemma above, we can obtain a sufficient condition for a numerical solution to converge to the exact solution of the ODE. This observation is specified in the next corollary.

\begin{corollary} \label{cor:conv_to_quad}
    Consider an array of dimension $N\geq1$ vectors $\{\pr{c^n}_j = \frac{1}{n} \sum_{l=0}^n \pr{\alpha^{(n)}_l}_j f(lh) \in \mathbb{R}^N \}_{n=1,2, \dots, \mathbb{N}}$, where $h = T/n$ and $f \colon [0,t] \to \mathbb{R}$ a Riemann integrable function.
    Assume that for each $j$, there exists a vector-valued degree $j-1$ polynomial function sequence $\{F^{(n)} \colon [0,1] \to \mathbb{R}^N \}_{n \in \mathbb{N}}$
    % where all of its elements interpolate the following $n-1$ points
    % $
    %     \pr{\frac{1}{n}, \pr{\alpha^{(n)}_1}_j}, \pr{\frac{2}{n}, \pr{\alpha^{(n)}_2}_j}, \dots, \pr{1 - \frac{1}{n}, \pr{\alpha^{(n)}_{n-1}}_j},
    % $ so that
    satisfying the following condion
    % and $f \colon [0,t] \to \mathbb{R}$ a Riemann integrable function. Assume $\sup_{n, l} \pr{\alpha^{(n)}_l}_j \leq C_j$. Denote the $j$-th Legendre polynomial as $P_j$.   
    % Now, consider a uniformly bounded function sequence $\{F^{(n)}\}_n$ such that $F \in C([0,1]; \mathbb{R}^N)$ for all $F \in \{F^{(n)}\}_{n \in \mathbb{N}}$. Further, assume that for all $F \in \{F^{(n)}\}_{n \in \mathbb{N}}$, the $j$-th component satisfies for all $j \in J$ and $l \in \{1, 2, \dots, n-1\}$
    \begin{align} \label{eq:F_condition}
        \pr{F \pr{\frac{l}{n}}}_j &= \pr{\alpha^{(n)}_l}_j
    \end{align}
    for all $l \in \{1, 2, \dots, n-1\}$.
    Further assume that $\left \|  \pr{F^{(n)}(\cdot)}_j - \sqrt{2j-1} P_{j-1} \pr{2\cdot-1} \right \|_{\sup([0,1])} \to 0$ holds as $n \to \infty$ for all $j \in J$. Then, 
    % \begin{align*}
    %     \pr{F^{(n)}(x)}_j \to \sqrt{2j-1} P_{j-1} \pr{2x-1} 
    % \end{align*}
    % uniformly on $x\in [0,1]$ as $n \to \infty$ for all $j \in J$, 
    \[
% \sup_{n=0,1,\dots, n}
\|c^n-c(t_n)\|\rightarrow 0,\qquad\text{ as } \quad n \rightarrow \infty.
    \]
    % \begin{align*}
    %     \pr{c^n}_j \to \frac{\sqrt{2j-1}}{t}\int_0^t P_{j-1}\pr{\frac{2s}{t}-1}f(s)ds
    % \end{align*}
    % as $h \to 0$ and $n \to \infty$ for all $j \in J$. 
    % Define the $j$-th component of the interpolating function for fixed $n$ as the polynomial with the lowest degree of order interpolating the following $n-1$ points
    % \begin{align} \label{inter_points}
    %     \pr{\frac{1}{n}, \alpha^{(n)}^j_1}, \pr{\frac{2}{n}, \alpha^{(n)}^j_2}, \dots, \pr{1 - \frac{1}{n}, \alpha^{(n)}^j_{n-1}}
    % \end{align}
    % as $F^{(n)}_j$. 
\end{corollary}

\begin{proof}
    Since the sequence $\{F^{(n)}\}$ is a polynomial sequence defined on a compact domain, with fixed degree of order, it is uniformly bounded. 
    Hence we let $F_j \leq M_j$ for all $F \in \{F^{(k)}\}$ and $\sup_{n, l} \pr{\alpha^{(n)}_l}_j \leq C_j$.
    Also, $|f| \leq M$ from the definition of Riemann integrable functions. Fix component index $j$. 
    
    Since $\pr{F^{(k)}\pr{\frac{l}{n}}}_j = \pr{\alpha^{(n)}_l}_j$ for $l \in \{1, 2, \dots, n-1\}$, we can write
    \begin{align*}
        \pr{c^n}_j &= \frac{1}{n} \sum_{l=1}^n \pr{F^{(n)}\pr{\frac{l}{n}}}_j f(lh) + \frac{1}{n} \pr{\pr{\alpha^{(n)}_0}_j f(0) + \pr{\alpha^{(n)}_n}_j f(nh) - \pr{F^{(n)}(1)}_j f(nh)}.
    \end{align*}
    Fix $\epsilon >0$. 
    Since the function sequence $F^{(n)}$ satisfies the conditions specified in Lemma \ref{lem:converge_condition}, we can find $N_1 \in \mathbb{N}$ such that for all $n \geq N_1$, $\left| \frac{1}{n} \sum_{l=1}^n \pr{F^{(n)}\pr{\frac{l}{n}}}_j f(lh) - \frac{\sqrt{2j-1}}{t}\int_0^t P_{j-1}\pr{\frac{2s}{t}-1}f(s)ds \right| < \epsilon / 2$. Choose $n \in \mathbb{N}$ so that $n \geq \max\{N_1, \frac{2M(2C_j+M_j)}{\epsilon}\}$. Constructing the following triangular inequality,
    \begin{align*}
        \left| \pr{c^n}_j - \frac{\sqrt{2j-1}}{t}\int_0^t P_{j-1}\pr{\frac{2s}{t}-1}f(s)ds \right| 
        & \leq \left| \frac{1}{n} \sum_{l=1}^n \pr{F^{(n)}\pr{\frac{l}{n}}}_j f(lh) - \frac{\sqrt{2j-1}}{t}\int_0^t P_{j-1}\pr{\frac{2s}{t}-1}f(s)ds \right| \\
        &+ \frac{1}{n} \left| \pr{\alpha^{(n)}_0}_j f(0) + \pr{\alpha^{(n)}_n}_j f(nh) \right| + \frac{1}{n} \left| \pr{F^{(n)}(1)}_j f(nh) \right|\\
        &\leq \frac{\epsilon}{2} + \frac{M}{n}\pr{2C_j+M_j} 
        \leq \epsilon
    \end{align*}
    holds for all $j \in \{1, 2, \dots, N\}$. 
    % where the convergence of $(1)$ follows from Lemma \ref{lem:converge_condition}.
\end{proof}

% Letting $G$ be the shifted Legendre Polynomial $P_j\pr{2x-1}$, the above lemma implies that the numerical solution of the form $\pr{I_n(f)}_j = \sum_{l=0}^n \alpha^{(n)}_l^j f(lh)$ will converge to the exact solution $\frac{2j-1}{t} P_{j-1}\pr{\frac{2s}{t}-1}f(s)ds$ given that the interpolating function converges to the shifted Legendre polynomial in $C([0,1], \|\cdot\|)$. 
The result of this corollary implies that instead of directly proving the convergence of the numerical solution $c^n$ to the exact solution, it would suffice to find a function sequence $\{F^{(n)}\}_{n \in \mathbb{N}}$ satisfying \eqref{eq:F_condition} that converges to the (scaled) shifted Legendre polynomial. 
A natural choice to construct such a sequence would be to interpolate the $n-1$ points using polynomials so that the function would satisfy \eqref{eq:F_condition}.
% Note that whether the starting index of the summation in $I_n(f)$ is 0 or 1 is not important since the contribution of a single point diminishes as $n \to \infty$. 
% The next task is to characterize the interpolating functions $F^j_n$. 
However, the degree of the interpolating polynomials in the sequence could diverge as $n \to \infty$, complicating the analysis of their limiting behavior. 
Surprisingly, the following lemma shows that for all discretization methods of interest, the sequence $\{F^{(n)}\}_{n \in \mathbb{N}}$ is a polynomial sequence of fixed degree. 
Note that if we interpolate $n+1$ points, i.e. including $\pr{0, \pr{\alpha^{(n)}_0}_j}, \pr{1, \pr{\alpha^{(n)}_n}_j}$ as interpolating points, the following lemma does not work. 

% {\color{red} How about changing this to proposition?}
\begin{lemma} \label{lem:degree_bound}
    Denote the $j$-th index of the numerical solution obtained by a given discretization method as $\pr{c^n}_j = \frac{1}{n} \sum_{l=0}^n \pr{\alpha^{(n)}_l}_j f^l$ where $h = t/n$. 
    Consider the following $n-1$ points 
    % Define the $j$-th component of the interpolating function for fixed $n$ as the polynomial with the lowest degree of order interpolating the following $n-1$ points
    \begin{align} \label{inter_points}
        \pr{\frac{1}{n}, \pr{\alpha^{(n)}_1}_j}, \pr{\frac{2}{n}, \pr{\alpha^{(n)}_2}_j}, \dots, \pr{1 - \frac{1}{n}, \pr{\alpha^{(n)}_{n-1}}_j}.
    \end{align}
    Then, for all $n \in \mathbb{N}$ and $j \in J$, there exists a degree $j-1$ polynomial $F^{(n)}_j$ that interpolates the above $n-1$ points obtained by 
    any discretization method introduced in Section~\ref{sec:problem setting}.
    % the forward Euler, backward Euler, bilinear, and approximated bilinear method. 
    Moreover, define $\{F^{(n)}\}_{n \in \mathbb{N}}$ as a vector-valued function sequence that for each $n$, the $j$-th component is $F^{(n)}_j$. 
    Then, given the eigenvalue decomposition $A=VDV^{-1}$, the sequence $\{F^{(n)}(x)\}_{n \in \mathbb{N}}$ pointwise converges to 
    \begin{align*}
        F(x) &= V \text{diag} \pr{ 1, 2x, 3x^2, \dots, Nx^{N-1}} V^{-1} e_1
    \end{align*}
    as $n \to \infty$ for all $x \in [0,1]$.
\end{lemma}

\begin{proof} 
    We start with the forward Euler case. Fix $n$. Let $p^{(n)}\colon [0,1] \to \mathbb{R}^N$ be some vector whose $j$-th component is a function interpolating the $n-1$ points in \eqref{inter_points}. Then, 
    referring to \eqref{eq:forward_numsol}, $p^{(n)}$ by construction satisfies
    \begin{align*}
        p^{(n)}\pr{\frac{l}{n}} 
    = \frac{n}{l}\prod_{k=l+1}^{n-1}\pr{ I - \frac{1}{k}A } A e_1
    \end{align*}
    for $l \in \{1, 2, ..., n-1\}$. Let $A = VDV^{-1}$ where $D \in \mathbb{R}^{N \times N}$ is the diagonal matrix with entries $(D)_{jj}=j$. Then, 
    \begin{align*}
        \prod_{k=l+1}^{n-1}\pr{ I - \frac{1}{k}A }
        &= \prod_{k=l+1}^{n-1} V \pr{ I - \frac{1}{k}D }V^{-1} 
        = V \pr{ \prod_{k=l+1}^{n-1} \pr{ I - \frac{1}{k}D } }V^{-1} \\
        &= V \pr{ \text{diag}\pr{ \prod_{k=l+1}^{n-1}\pr{ 1 - \frac{1}{k} }, \prod_{k=l+1}^{n-1}\pr{ 1 - \frac{2}{k} }, \dots, \prod_{k=l+1}^{n-1}\pr{ 1 - \frac{j}{k} }, \dots, \prod_{k=l+1}^{n-1}\pr{ 1 - \frac{N}{k} }  } }V^{-1}
    \end{align*}
    Now, since we are interested in the limiting behavior of $n \to \infty$, we can assume that $n$ is considerably larger than $N$. 
    Then we can cancel out terms in the denominator and the numerator to calculate the $i$-th term in the diagonal matrix, 
    \begin{align*}
        \frac{1}{l}\prod_{k=l+1}^{n-1}\pr{ 1- \frac{i}{k} } 
        &=  \frac{1}{l} \frac{1}{\prod_{k=l+1}^{n-1} k} \prod_{k=l+1}^{n-1}\pr{ k - i }
        % &= \frac{(l+1 -i)(l+2 -i)\dots(n-1-i)}{l(l+1) \dots (n-1)} \\
        % &= \frac{(l+1-i)\dots(l-1)}{(n-i) \dots (n-2)(n-1)} \\
        = \frac{\prod_{k=1}^{i-1}\pr{ l - k }}{\prod_{k=1}^{i}\pr{ n - k }}
    \end{align*}
     where $\prod_{k=1}^{0}\pr{ l - k } = 1$. Therefore we arrive at
     \begin{align*}
    p^{(n)}\pr{\frac{l}{n}} 
        = n V \text{diag}\pr{ \frac{1}{n-1}, \frac{ \pr{ l-1}}{(n-1)(n-2)}, \dots, \frac{\prod_{k=1}^{N-2}\pr{ l - k } }{{\prod_{k=1}^{N-1}\pr{ n - k }}}, \frac{\prod_{k=1}^{N-1}\pr{ l - k } }{{\prod_{k=1}^{N}\pr{ n - k }}}} D V^{-1} e_1
     \end{align*}
     for $l \in \{1, 2, \dots, n-1\}$. 
     Now we change variables and let $l = nx$. Then, we can define a vector function $F^{(n)}$ with the $j$-th component
     \begin{align*}
         F_j^{(n)}(x) &= n e_j^{t} V \text{diag}\pr{ \frac{1}{n-1}, \frac{ \pr{ nx-1}}{(n-1)(n-2)}, \dots, \frac{\prod_{k=1}^{N-2}\pr{ nx - k } }{{\prod_{k=1}^{N-1}\pr{ n - k }}}, \frac{\prod_{k=1}^{N-1}\pr{ nx - k } }{{\prod_{k=1}^{N}\pr{ n - k }}}} D V^{-1} e_1.
     \end{align*}
    such that $F^{(n)}\pr{\frac{l}{n}} = p^{(n)}\pr{\frac{l}{n}}$ for all $l = \{1,2, \dots, n-1\}$, for all $x \in [0,1]$.      
     Notice in the above expression of $F_j^{(n)}(x)$ that the $i$-th term in the diagonal matrix is a $i-1$ degree polynomial of $x$. Since $V$ and $V^{-1}$ are both lower triangular, we conclude that $F_j^{(n)}$ is a $j-1$ degree polynomial interpolating the $n-1$ points of interest. 
    Moreover, for any fixed $x \in [0,1]$, taking the limit $n \to \infty$ yields
    \begin{align} \label{eqref:F^n_limit}
         \lim_{n \to \infty} F^{(n)}(x) &= V \text{diag} \pr{ 1, 2x, 3x^2, \dots, Nx^{N-1}} V^{-1} e_1.
     \end{align}
     This concludes the proof for the forward Euler method. 

     We use a similar approach for other methods. For the backward Euler method, referring to \eqref{eq:back_numsol}, the interpolating function satisfies
     \begin{align*}
         p^{(n)}\pr{\frac{l}{n}} &= \frac{n}{l}\prod_{k=l}^n \pr{I + \frac{1}{k}A}^{-1}Ae_1
     \end{align*}
     for $l \in \{1, 2, \dots, n-1\}$. Then the product term in the RHS is equal to
     \begin{align*}
         \prod_{k=l}^{n}\pr{ I + \frac{1}{k}A }^{-1}
        &= \prod_{k=l+1}^{n} V \pr{ I + \frac{1}{k}D }^{-1} V^{-1}
        = V \pr{ \prod_{k=l}^{n} \pr{ I + \frac{1}{k}D } }^{-1} V^{-1} \\
        &= V \pr{ \text{diag}\pr{ \prod_{k=l}^{n}\pr{ 1 + \frac{1}{k} }^{-1}, \prod_{k=l}^{n}\pr{ 1 + \frac{2}{k} }^{-1}, \dots, \prod_{k=l}^{n}\pr{ 1 + \frac{j}{k} }^{-1}  } } V^{-1}
     \end{align*}
     and by canceling out terms assuming $n$ is large,
     \begin{align*}
         \frac{1}{l}\prod_{k=l}^{n}\pr{ 1 + \frac{i}{k} }^{-1} 
        &=  \frac{1}{l}  \frac{\prod_{k=l}^{n} k}{\prod_{k=l}^{n} \pr{ k + i }}
        % &= \frac{(l+1)(l+2)\cdots(n-1)n}{(l+i)(l+i+1)\cdots(n+i-1)(n+i)} \\
        % &= \frac{(l+1)\cdots(l+i-1)}{(n+1) \cdots (n+i-1)(n+i)} \\
        = \frac{\prod_{k=1}^{i-1}\pr{ l + k }}{\prod_{k=1}^{i}\pr{ n + k }}.
     \end{align*}
    Similar as in the forward Euler case, we can define the interpolating polynomial $F_j^{(n)}$ as 
     \begin{align*}
          F_j^{(n)}(x) &= n e_j^{t} V \text{diag}\pr{ \frac{1}{n+1}, \frac{ \pr{ nx+1}}{(n+1)(n+2)}, \dots, \frac{\prod_{k=1}^{N-2}\pr{ nx + k } }{{\prod_{k=1}^{N-1}\pr{ n + k }}}, \frac{\prod_{k=1}^{N-1}\pr{ nx + k } }{{\prod_{k=1}^{N}\pr{ n + k }}}} D V^{-1} e_1.
     \end{align*}
     By the same reasoning as for the forward Euler case, we conclude that $F_j^{(n)}$ is a $j-1$ degree polynomial. Moreover, it converges to \eqref{eqref:F^n_limit} pointwise as $n \to \infty$ for $x \in [0,1]$ as $n \to \infty$.

     For the bilinear method, referring to \eqref{eq:bilin_numsol}, the interpolating function satisfies
     \begin{align*}
         p^{(n)}\pr{\frac{l}{n}} &= \frac{n}{2l} \pr{I+\frac{1}{2n}A}^{-1} 
         \pr{
         \prod_{k=l}^{n-1} 
         \pr{I - \frac{1}{2k}A}\pr{I + \frac{1}{2k}A}^{-1}
         + \prod_{k=l+1}^{n-1} 
         \pr{I - \frac{1}{2k}A}\pr{I + \frac{1}{2k}A}^{-1}
         }Ae_1 \\
         &= \frac{n}{2l} \pr{I+\frac{1}{2n}A}^{-1} 
        \pr{I + \pr{I-\frac{1}{2l}A}\pr{I+\frac{1}{2l}A}^{-1}}
         \prod_{k=l+1}^{n-1} 
         \pr{I - \frac{1}{2k}A}\pr{I + \frac{1}{2k}A}^{-1}
         Ae_1
     \end{align*}
     for $l \in \{1, 2, \dots, n-1\}$. 
     For the RHS, the last term simplifies to 
    \begin{align*}
        \prod_{k=l+1}^{n-1} 
         \pr{I - \frac{1}{2k}A}\pr{I + \frac{1}{2k}A}^{-1} 
         &= V \prod_{k=l+1}^{n-1}  \pr{I - \frac{1}{2k}D}\pr{I + \frac{1}{2k}D}^{-1}  V^{-1} \\ 
         &= V \pr{ \text{diag}\pr{ \prod_{k=l+1}^{n-1}\pr{\frac{k-1/2}{k+1/2} }, \prod_{k=l+1}^{n-1}\pr{\frac{k-1}{k+1} }, \dots, \prod_{k=l+1}^{n-1} \pr{ \frac{k-j/2}{k+j/2} }}}V^{-1}
    \end{align*}
    and the prefix terms simplify to 
    \begin{align*}
        &\frac{n}{2l} \pr{I+\frac{1}{2n}A}^{-1} 
        \pr{I + \pr{I-\frac{1}{2l}A}\pr{I+\frac{1}{2l}A}^{-1}}\\
        &= \frac{n}{2l} V \pr{I+\frac{1}{2n}D}^{-1} 
        \pr{I + \pr{I-\frac{1}{2l}D}\pr{I+\frac{1}{2l}D}^{-1}} V^{-1} \\
        &= V \pr{\text{diag} \pr{ 
        \pr{\frac{n^2}{(n+1/2)(l+1/2)}}, \pr{\frac{n^2}{(n+1)(l+1)}}, \dots, \pr{\frac{n^2}{(n+1/2j)(l+1/2j)}}}} V^{-1}.
    \end{align*}
    Combining these two terms and simplifying the denominators and numerators, we obtain 
     % \begin{align*}
     %     \frac{1}{l}\prod_{k=l}^{n}\pr{ 1 + \frac{i}{k} }^{-1} 
     %    &=  \frac{1}{l}  \frac{\prod_{k=l}^{n} k}{\prod_{k=l}^{n} \pr{ k + i }} \\
     %    &= \frac{(l+1)(l+2)\cdots(n-1)n}{(l+i)(l+i+1)\cdots(n+i-1)(n+i)} \\
     %    &= \frac{(l+1)\cdots(l+i-1)}{(n+1) \cdots (n+i-1)(n+i)} \\
     %    &= \frac{\prod_{k=1}^{i-1}\pr{ l + k }}{\prod_{k=1}^{i}\pr{ n + k }}
     % \end{align*}
    \begin{align*}
        p^{(n)}\pr{\frac{l}{n}} 
        = V \text{diag}\pr{ \frac{4n^2}{(2n-1)(2n+1)}, \frac{ ln}{(n-1)(n+1)}, \dots, \frac{\prod_{k=1}^{N}\pr{ l - N/2 + k } }{{\prod_{k=1}^{N}\pr{ n - N/2 + k - 1}}} 
        \cdot \frac{n^2}{(l+j/2)(n+j/2)} } D V^{-1} e_1
    \end{align*}
    % \begin{align*}
    %     p^{(n)}\pr{\frac{l}{n}} 
    %     = P^{-1} \text{diag}\pr{ \frac{4n^2}{(2n-1)(2n+1)}, \frac{ l}{n(n-1)(n+1)}, \dots, \frac{\prod_{k=1}^{j}\pr{ l - j/2 + k } }{{\prod_{k=1}^{j}\pr{ n - j/2 + k - 1}}} 
    %     \cdot \frac{1}{(l+j/2)(n+j/2)} } D P e_1
    % \end{align*}

     Define $F_j^{(n)}$ as 
     \begin{align*}
        F_j^{(n)}(x) 
        = e_j^tV \text{diag}\pr{ \frac{4n^2 }{(2n-1)(2n+1)}, 
        \frac{n^2x}{(n-1)(n+1)},
        \dots,
        \frac{\prod_{k=1}^{N-1}\pr{ nx - N/2 + k } }{{\prod_{k=1}^{N}\pr{ n - N/2 + k - 1}}} \cdot
        \frac{n^2}{n+N/2}} D V^{-1} e_1.
     \end{align*}
     By the same reasoning as for the forward Euler case, we conclude that $F_j^{(n)}$ is a $j-1$ degree polynomial. Moreover, it converges to \eqref{eqref:F^n_limit} pointwise as $n \to \infty$ for $x \in [0,1]$ as $n \to \infty$.

    For the approximated bilinear method, 
    referring to \eqref{eq:apbil_numsol}, the interpolating function satisfies
     \begin{align*}
         p^{(n)}\pr{\frac{l}{n}} &= \frac{n}{l}
         \prod_{k=l+1}^n 
         \pr{I - \frac{1}{2k}A}
         \prod_{k=l}^n 
         \pr{I + \frac{1}{2k}A}^{-1}
         Ae_1 \\
         &= \frac{n}{l}
         \pr{ I + \frac{1}{2l}A}^{-1}
         \prod_{k=l+1}^n 
         \pr{I - \frac{1}{2k}A} \pr{I + \frac{1}{2k}A}^{-1}
         Ae_1
     \end{align*}
     for $l \in \{1, 2, \dots, n-1\}$. 
     Assuming $n$ is large, the product term in RHS simplifies as
     \begin{align*}
         \prod_{k=l+1}^n 
         \pr{I - \frac{1}{2k}A} \pr{I + \frac{1}{2k}A}^{-1}
        &= \prod_{k=l+1}^n V 
         \pr{I - \frac{1}{2k}D} \pr{I + \frac{1}{2k}D}^{-1} V^{-1} \\
        &= V \pr{ \text{diag}
        \pr{ \prod_{k=l+1}^{n}\pr{ \frac{k-1/2}{k+1/2} }, \prod_{k=l+1}^{n}\pr{ \frac{k-1}{k+1} }, \dots, \prod_{k=l+1}^{n}\pr{ \frac{k - j/2}{k + j/2}}}}
        V^{-1}
     \end{align*}
     and the prefix term simplifies to
     \begin{align*}
         \frac{1}{l}\pr{ I + \frac{1}{2l}A}^{-1} 
         &= \frac{1}{l} P^{-1} \pr{ I + \frac{1}{2l}D}^{-1} P 
         = V \pr{ \text{diag} \pr{\frac{1}{l+1/2}}, \pr{\frac{1}{l+3/2}} \dots,
         \pr{\frac{1}{l+j/2}} }V^{-1}.
     \end{align*}
     Assuming that $n$ is large and canceling out terms, we obtain
    \begin{align*}
        p^{(n)}\pr{\frac{l}{n}} 
        = V \text{diag}
        \pr{ \frac{2n}{2n+1}, \frac{ln}{n(n+1)}, \dots, \frac{\prod_{k=1}^{N}\pr{ l - N/2 + k } }{{\prod_{k=1}^{N}\pr{ n - N/2 + k}}} 
        \cdot \frac{n}{l+N/2} } D V^{-1} e_1
    \end{align*}
    
    Define $F_j^{(n)}$ as 
     \begin{align*}
          F_j^{(n)}\pr{x} 
        = e_j^t V \text{diag}
        \pr{ \frac{2n}{2n+1}, \frac{n^2x}{n(n+1)}, \dots, \frac{n\prod_{k=1}^{j-1}\pr{ nx - j/2 + k } }{{\prod_{k=1}^{j}\pr{ n - j/2 + k}}}} D V^{-1} e_1
     \end{align*}
     By the same reasoning as in the forward Euler case, we conclude that $F_j^{(n)}$ is a $j-1$ degree polynomial. Moreover, it converges to \eqref{eqref:F^n_limit} pointwise as $n \to \infty$ for $x \in [0,1]$ as $n \to \infty$.

    For zero-order hold, referring to \eqref{eq:zoh_numsol}, the interpolating function satisfies
    \begin{align*}
        p^{(n)} \pr{\frac{l}{n}} &= n \pr{e^{A\log\pr{\frac{l+1}{n}}} - e^{A\log\pr{\frac{l}{n}}}} e_1
    \end{align*}
    for $l \in \{1,2, \dots, n-1\}$. Diagonalize $A$ by $A = VDV^{-1}$ and 
    \begin{align*}
        p^{(n)} \pr{\frac{l}{n}} &= n V \pr{e^{D\log\pr{\frac{l+1}{n}}} - e^{D\log\pr{\frac{l}{n}}}} V^{-1} e_1 \\
        &= V \text{diag}
        \pr{1, \frac{1}{n} \pr{\pr{l+1}^2 - l^2}, \dots, 
        \frac{1}{n^{j-1}} \pr{\pr{l+1}^j-l^j}, \dots,
        \frac{1}{n^{N-1}} \pr{\pr{l+1}^N-l^N}
        } V^{-1} e_1.
        % \pr{\pr{\frac{l+1}{n}}^j - \pr{\frac{l}{n}}^j}e_1
    \end{align*}
    Similarly as in the other cases, we can define $F_j^{(n)}$ as 
    \begin{align*}
        F_j^{(n)} (x)
        &= e_j^t V \text{diag} \pr{
        1, 2x + \frac{1}{n}, 
        \frac{1}{n^{j-1}} \sum_{k=0}^{j-1} \binom{j}{k} \pr{nx}^k, \dots,
        \frac{1}{n^{N-1}} \sum_{k=0}^{N-1} \binom{N}{k} \pr{nx}^k
        }V^{-1}e_1
    \end{align*}
    which pointwise converges to \eqref{eqref:F^n_limit}.
\end{proof}

Combining the results, we can prove Theorem \ref{thm:conv_riemann}.
% the following theorem shows that all the discretization methods introduced in Section \ref{subsec:disc_methods} are convergent for every Riemann integrable $f$. 
% The following theorem generalizes Theorem \ref{thm:conv_f_cont} by allowing $f$ to be Riemann integrable and proving convergence for the approximated bilinear method, successfully filling in all the gaps introduced in Section \ref{sec:error_analysis}.



\begin{proof} [Proof of Theorem \ref{thm:conv_riemann}]
    Given a discretization method, we can express the numerical solution as $c^n = \frac{1}{n} \sum_{l=0}^n \alpha^{(n)}_l f^l$ by Lemma \ref{lem:numsol_linear}.
    Then, define the function sequence $\{F^{(n)}\}_{n \in \mathbb{N}}$ as in Lemma \ref{lem:degree_bound}. Then for all $n \in \mathbb{N}$, the $j$-th component of $F^{(n)}\colon [0,1] \to \mathbb{R}^N$ is a $j-1$ degree polynomial with pointwise limit $F(x)$ for all $x \in [0,1]$. This implies that all the coefficients of the polynomials in the sequence converge to the coefficients of $F$, and hence we can conclude that $\{F^{(n)}\}_{n \in \mathbb{N}}$ converges to $F$ uniformly. 
    Note that we also have that all the coefficients $\alpha^{(n)}_l$ of $c^n$ are uniformly bounded, since the limit is well-defined.

    Now it suffices to show that the $j$-th component of the limit function $F(x)\colon [0,1] \to \mathbb{R}^N$ is equal to $\sqrt{2j-1}$ times the $j-1$-th shifted Legendre polynomial, $\sqrt{2j-1}P_{j-1}(2x-1)$. 
    Once this is shown, we can apply the result of Corollary \ref{cor:conv_to_quad} to conclude the proof.
    
    Recall that the exact form of $F$ is:
    \begin{align*}
        F(x) &= V \text{diag}\pr{1, 2x, 3x^2, \dots, jx^{j-1}, \dots, Nx^{N-1}} V^{-1} e_1.
    \end{align*}

    Differentiating both sides with respect to $x$, we get
    \begin{align*}
        F'(x) &= V \text{diag}\pr{0, 2, 6x, \dots, j(j-1)x^{j-2}, \dots, N(N-1)x^{N-2}} V^{-1} e_1.
    \end{align*}

    Combining these two equations, we obtain the following differential equation that holds for all $x \in [0,1]$:
    \begin{align} \label{eq:F_recur}
        xF'(x) &= (A - I) F(x)
    \end{align}

    Since we know the exact form of $A$, we can derive a recurrence relation for $F$ for arbitrary dimension $N$. 
    Rewriting the $j$-th component $F$ as
    \begin{align*}
        % F(x) = \begin{bmatrix}
        %     f_0(x) \\
        %     \sqrt{3} f_1(x)\\
        %     \sqrt{5} f_2(x)\\
        %     \dots
        % \end{bmatrix},
        F_j(x) &= \sqrt{2j-1}f_{j-1}(x),
    \end{align*}
    we obtain the recurrence relation 
    \begin{align} \label{eq:f_recurrence}
        xf'_j(x) = jf_j(x) + \sum_{l=0}^{j-1} (2l+1)f_{l}(x).
        % jf_j(x) + (2j-1)f_{j-1}(x) + (2j-3)f_{j-2}(x) + \dots
    \end{align}
    Notice that \eqref{eq:f_recurrence} is exactly the recurrence relation satisfied the $j$-th shifted Legendre polynomial $\tilde{P}_j$. 
    Matching the initial condition $F(1) = Ae_1 = B$, we have $f_j(1) = 1$ for all $j \in J$. Then by induction, we can prove that $f_j(x) = \tilde{P}_j(x) = P_j(2x-1)$.
    Finally, utilizing the uniqueness of the solution for the IVP defined with ODE \eqref{eq:F_recur} and initial condition at $x=1$, we arrive at the conclusion:
    \begin{align*}
        % F(x) &= \begin{bmatrix}
        %     P_0(2x-1) \\
        %     \sqrt{3} P_1(2x-1) \\
        %     \sqrt{5} P_2(2x-1) \\
        %     \dots
        % \end{bmatrix}.
        F_j(x) &= \sqrt{2j-1}P_{j-1}(2x-1).
    \end{align*}
    % Given a discretization method and number of mesh points $n$, denote the obtained numerical solution as $c_n = I_n(f) = \frac{1}{n} \sum_{l=0}^n \alpha^{(n)}_l f(lh)$. 
\end{proof}


\section{Convergence rate analysis}
In this section, we analyze the convergence rates of the numerical solutions obtained by the discretization methods introduced in Section \ref{sec:problem setting}. While \emph{convergence} was proven for all Riemann integrable input functions $f$ in Section \ref{subsec:riemann_f}, the analysis of \emph{convergence rates} will require additional regularity assumptions.

% This is important because this could give a worst-case bound on the order of accuracy we have given the stepsize. This provides for a reasonable tradeoff between accuracy and compute. 

% \subsection{Theoretical analysis} \label{subsec:conv_rate_theory}
% By imposing stronger conditions on input function $f$, we can obtain certain convergence rates.

\begin{theorem} [Convergence rates of numerical solutions] \label{thm:conv_rate}
    Consider the setup in Theorem \ref{thm:conv_riemann}. Assume further that $f$ is of bounded variation on $[0,T]$. Then, we can obtain $\mathcal{O}(1/n)$ convergence rate for forward Euler, backward Euler, bilinear, approximated bilinear, and zero-order hold.     
    Further assume that $f \in C^2([0,T])$. Then, we attain $\mathcal{O}(1/n^2)$ convergence rate for the bilinear method.
\end{theorem}

\begin{proof}
Define $T(n) = \frac{1}{n}\sum_{l=0}^{n}
P_{j-1}\pr{\frac{2l}{n}-1} f(lh) - \frac{1}{2n} \pr{P_{j-1}(-1) f(0) + P_{j-1}(1) f(n)}$ as the result of applying the composite trapezoidal rule to the exact solution. Then we can construct a triangle inequality as
\begin{align} \label{ineq:conv_rate_start}
    \left| \pr{c^n}_j - \frac{\sqrt{2j-1}}{t}\int_0^t P_{j-1}\pr{\frac{2s}{t}-1}f(s)ds \right| 
    % & \leq \left| \frac{1}{n} \sum_{l=1}^n \pr{F^{(n)}\pr{\frac{l}{n}}}_j f(lh) - \frac{\sqrt{2j-1}}{t}\int_0^t P_{j-1}\pr{\frac{2s}{t}-1}f(s)ds \right| \\
    % &+ \frac{1}{n} \left| \pr{\alpha^{(n)}_0}_j f(0) + \pr{\alpha^{(n)}_n}_j f(nh) \right| + \frac{1}{n} \left| \pr{F^{(n)}(1)}_j f(nh) \right|\\
    &\leq 
    \underbrace{\left|  \pr{c^n}_j 
    -  \sqrt{2j-1} T(n) \right|}_{(1)}\\
    &+ \underbrace{\sqrt{2j-1} \left| T(n)
    - \frac{1}{t}\int_0^t P_{j-1}\pr{\frac{2s}{t}-1}f(s)ds \right|}_{(2)}. \nonumber
    % &+ \underbrace{\frac{M}{n}\pr{2C_j+M_j}}_{(3)}
\end{align}
$(1)$ corresponds to the error between the numerical solution and the Riemann sum (obtained by applying the trapezoidal rule).
For $(2)$, we prove the following lemma.
\begin{lemma}
    Assume that a real-valued function $f$ is of bounded variation on the interval $[0,t]$ with fixed $t$. Then, 
    \begin{align*}
        \abs{\frac{1}{t} \int_0^t f(s)ds - \frac{1}{n} \sum_{i=1}^{n}f\pr{\frac{it}{n}}} \leq \mathcal{O}(1/n).
    \end{align*}
\end{lemma}

\begin{proof}
    Denote $h = t/n$ and $V_a^b(f)$ be the total variation of $f$ on a closed interval $[a,b]$. Then,
    \begin{align*}
        \abs{
        \frac{1}{nh}\int_0^{nh}f(s)ds - \frac{1}{n} \sum_{i=1}^{n}f\pr{ih}
        }
        &\leq \frac{1}{nh} \sum_{i=1}^n 
        \int_{(i-1)h}^{ih} \abs{f(s) - f(ih)} ds\\
        &\leq \frac{1}{n} \sum_{i=1}^n \pr{ \sup_{x \in [(i-1)h, ih]} f(x) - \inf_{x \in [(i-1)h, ih]} f(x) } \\
        &\leq \frac{1}{n} \sum_{i=1}^n V_{(i-1)h}^{ih}(f)
        \leq \frac{V_0^t (f)}{n}.
        % &\leq \frac{1}{n} \sum_{i=1}^n 
        % \abs{ \max_{x \in [(i-1)h, ih]} f(x) - \min_{x \in [(i-1)h, ih]} f(x) } \\
        % &\leq \frac{1}{n} \sum_{i=1}^{n_\mathcal{P}} 
        % \abs{ f(x_{i+1}) - f(x_i) } 
        % \leq \frac{1}{n} V(f)
    \end{align*}
    % where $\mathcal{P}$ is a partition of interval $[0,t]$ that includes all $ih$, $\text{argmax}_{x \in [(i-1)h, ih]}f(x)$, and $\text{argmin}_{x \in [(i-1)h, ih]}f(x)$ for all $i \in \{1, 2, \dots, n\}$, and $x_i$'s are its partition points. 
\end{proof}

Arranging some terms and applying the lemma above, one can conclude that $(2)$ has asymptotic rate of $O(1/n)$. 
For the asymptotic rate of $(1)$, 
recall that for $\pr{c^n}_j = \frac{1}{n} \sum_{l=0}^n \pr{\alpha_l^{(n)}}_j f^l$, the interpolating function $F^{(n)}_j$ was defined such that $F^{(n)}_j \pr{\frac{l}{n}} = \pr{\alpha^{(n)}_l}_j$ for $l = \{1, \dots, n-1\}$ (endpoints are excluded). It was proved in Lemma 
\ref{lem:degree_bound} and Theorem \ref{thm:conv_riemann} that
% that $F_j^{(n)}(x)$ converges to 
\begin{align*}
     \lim_{n \to \infty} F^{(n)}(x) &= F(x) %\\&
     = V \text{diag} \pr{ 1, 2x, 3x^2, \dots, Nx^{N-1}} V^{-1} e_1,
 \end{align*}
and that the $j$-th component of $F(x)$ is the scaled-shifted $(j-1)$-th Legendre polynomial, i.e., 
\begin{align*}
    F_j(x) = \sqrt{2j-1}P_{j-1}(2x-1).
\end{align*}
% We first calculate the rate for the forward Euler method (the same approach could be used for backward Euler, approximate bilinear, and zero-order hold). 
Rewriting $(1)$,
\begin{align*}
    \left|\pr{c^n}_j - \sqrt{2j-1} T(n) \right| 
    &= \left| \frac{1}{n} \sum_{l=0}^n \pr{\alpha_l^{(n)}}_j f^l - \sqrt{2j-1} T(n)
    \right| \\
    &\leq \frac{1}{n} \sum_{l=1}^{n-1} \left| F^{(n)}_j\pr{\frac{l}{n}}f^l - \sqrt{2j-1} P_{j-1}\pr{\frac{2l}{n}-1}f^l \right| \\
    &\phantom{=}+ \frac{1}{n} \left| \pr{\alpha^{(n)}_0}_j f^0 + \pr{\alpha^{(n)}_n}_j f^1 - \frac{\sqrt{2j-1}}{2}\pr{P_{j-1}(-1)f^0 + P_{j-1}(1)f^1} \right| \\
    &= \frac{M}{n} \sum_{l=1}^{n-1} \left| F^{(n)}_j\pr{\frac{l}{n}} - F_j\pr{\frac{l}{n}} \right| + \frac{K_n}{n} \\
    &\leq M \underbrace{\| F^{(n)}_j - F_j \|_{\sup}}_{(\star)} + \frac{K_n}{n}
\end{align*}
where $M$ is upper bound for $f$. 
Note that $K_n$ is uniformly bounded, so the second part automatically is of $\mathcal{O}(1/n)$. 
For the first term, note that $(\star)$ differs depending on which discretization method we are considering. 
Starting with the forward Euler method, recall that from Lemma~\ref{lem:degree_bound} that the interpolating function $F_j$ is defined as 
 \begin{align*}
     F_j^{(n)}(x) &= e_j^{t} V \text{diag}\pr{ \frac{n}{n-1}, \frac{ n \pr{ nx-1}}{(n-1)(n-2)}, \dots, \frac{n\prod_{k=1}^{N-2}\pr{ nx - k } }{{\prod_{k=1}^{N-1}\pr{ n - k }}}, \frac{n\prod_{k=1}^{N-1}\pr{ nx - k } }{{\prod_{k=1}^{N}\pr{ n - k }}}} D V^{-1} e_1.
 \end{align*}
for $x \in [0,1]$. 
% For all fixed $j$ and $x$, this is a rational function with respect to $n$, and ($j-1$ order) polynomial of $x$. 
Observing the diagonal components, 
for every $x\in[0,1]$, we can write the component of $F_j(x)-F_j^{(n)}(x)$ as $\frac{p_x(n)}{q(n)}$ with some polynomials $p_x$, $q$. Note $F_j(x)$ is constant with respect to $n$. Taking closer look at the numerators of $F_j^{(n)}(x)$, we can check the coefficients of $p_x(n)$ are polynomials with respect to $x$. Denote the coefficient of leading term as $a(x)$. 
Since $\lim_{n\to\infty} \frac{p_x(n)}{q(n)}=0$, we have $\lim_{n\to\infty} n\frac{p_x(n)}{q(n)}=a(x) \le \max_{x\in[0,1]} \abs{a(x)}$. As $F_j(x)-F_j^{(n)}(x)$ is finite dimensional, we conclude $\|F_j-F_j^{(n)}\|_{\sup} = \mathcal{O}(1/n)$ for all $j \in \{1, \dots, N\}$. 

% will have the leading term as $g_j(x)/n$, where $g(x)$ is some polynomial of $x$. Since $x\in[0,1]$, $g_j$ can always be bounded, and we conclude that $(\star)$ is $O(1/n)$. 
The same argument holds with the backward Euler method, approximated bilinear method, and zero-order hold, with the only difference in computing $(\star)$.
From Lemma~\ref{lem:degree_bound} we have 
% , all other steps are identical and we only need to check the asymptotic rate of $(\star)$. Recall $F_j^{(n)}$ (computed in Lemma~\ref{lem:degree_bound}) for each of these methods are:
\begin{align*}
    (F_j^{(n)})_{\text{backward}} (x)
    &= n e_j^{t} V \text{diag}\pr{ \frac{1}{n+1}, \frac{ \pr{ nx+1}}{(n+1)(n+2)}, \dots, \frac{\prod_{k=1}^{N-2}\pr{ nx + k } }{{\prod_{k=1}^{N-1}\pr{ n + k }}}, \frac{\prod_{k=1}^{N-1}\pr{ nx + k } }{{\prod_{k=1}^{N}\pr{ n + k }}}} D V^{-1} e_1 \\
    (F_j^{(n)})_{\text{approx bilin}} (x)
    &= e_j^t V \text{diag}
    \pr{ \frac{2n}{2n+1}, \frac{n^2x}{n(n+1)}, \dots, \frac{n\prod_{k=1}^{j-1}\pr{ nx - j/2 + k } }{{\prod_{k=1}^{j}\pr{ n - j/2 + k}}}} D V^{-1} e_1 \\
    (F_j^{(n)})_{\text{zoh}} (x)
    &= e_j^t V \text{diag} \pr{
    1, 2x + \frac{1}{n}, 
    \frac{1}{n^{j-1}} \sum_{k=0}^{j-1} \binom{j}{k} \pr{nx}^k, \dots,
    \frac{1}{n^{N-1}} \sum_{k=0}^{N-1} \binom{N}{k} \pr{nx}^k
    }V^{-1}e_1.
\end{align*}
Using the same argument, we conclude that the backward Euler, approximate bilinear and zero-order hold methods achieve $1/n$ convergence rate.

For the bilinear method, we now assume that $f \in C^2([0,T])$. Starting from \eqref{ineq:conv_rate_start}, we first know that $(2)$ is of $\mathcal{O}(1/n^2)$ by classical quadrature results for smooth $f$. For the asymptotic rate of $(1)$, we have to consider the asymptotic rate of both $(\star)$ and $K_n$. 
% we construct the same inequality
% \begin{align*}
%     \left|\pr{c^n}_j - \sqrt{2j-1} T(n) \right| 
%     &= \left| \frac{1}{n} \sum_{l=0}^n \pr{\alpha_l^{(n)}}_j f^l - \sqrt{2j-1} T(n)
%     \right| \\
%     &\leq \frac{1}{n} \sum_{l=1}^{n-1} \left| F^{(n)}_j\pr{\frac{l}{n}}f^l - \sqrt{2j-1} P_{j-1}\pr{\frac{2l}{n}-1}f^l \right| \\
%     &+ \frac{1}{n} \left| \pr{\alpha^{(n)}_0}_j f^0 + \pr{\alpha^{(n)}_1}_j f^1 - \frac{\sqrt{2j-1}}{2}\pr{P_{j-1}(-1)f^0 + P_{j-1}(1)f^1} \right| \\
%     &= \frac{M}{n} \sum_{l=1}^{n-1} \left| F^{(n)}_j\pr{\frac{l}{n}} - F_j\pr{\frac{l}{n}} \right| + \frac{1}{n} \left| \pr{\alpha^{(n)}_0}_j f^0 + \pr{\alpha^{(n)}_1}_j f^1 - \frac{\sqrt{2j-1}}{2}\pr{(-1)^{j-1}f^0 + f^1} \right| \\
%     &\leq M \sum_{l=1}^{n-1} \underbrace{\| F^{(n)}_j - F_j \|_{\sup}}_{(\star\star)} + \frac{1}{n} \left| \pr{\alpha^{(n)}_0}_j f^0 + \pr{\alpha^{(n)}_1}_j f^1 - \frac{\sqrt{2j-1}}{2}\pr{(-1)^{j-1}f^0 + f^1} \right|
% \end{align*}
The convergence rate for $(\star)$ could be obtained in a similar manner. 
From Lemma~\ref{lem:degree_bound}, $F_j^{(n)}$ obtained by applying the bilinear method is
 \begin{align*}
    F_j^{(n)}(x) 
    = e_j^t V \text{diag}\pr{ \frac{n^2 }{(n-1/2)(n+1/2)}, 
    \frac{n^2x}{(n-1)(n+1)},
    \dots,
    \frac{\prod_{k=1}^{N-1}\pr{ nx - N/2 + k } }{{\prod_{k=1}^{N}\pr{ n - N/2 + k - 1}}} \cdot
    \frac{n^2}{n+N/2}} D V^{-1} e_1.
 \end{align*}
for $x \in [0,1]$. 
Observe that denominator of the $j$-th term in the diagonal matrix is $\prod_{k=1}^{j}\pr{n - j/2 + k -1} \cdot (n + j/2) = (n-j/2)(n-j/2+1)\cdots(n+j/2-1)(n-j/2) = (n^2 - (j/2)^2)(n^2 - (j/2 -1)^2) \cdots$, where the last term is $n$ if $j$ is even and $(n^2 - (1/2)^2)$ if $j$ is odd. Either way, the subleading term (w.r.t. $n$) is 2 orders less than the leading term. 
Then, as before, writing a component of $F_j(x) - F_j^{(n)}$ as $\frac{p_x(n)}{q(n)}$, we have $\lim_{n\to\infty} n^2 \frac{p_x(n)}{q(n)} = a(x) \leq \max_{x\in [0,1]}|a(x)|$. Hence we conclude $\|F_j-F_j^{(n)}\|_{\sup} = \mathcal{O}(1/n^2)$. 

% For each term in the diagonal matrix, one can find that the subleading term (w.r.t. $n$) of the denominator is 2 orders less than the leading term (due to the symmetric multiplication).
% Taylor expanding w.r.t. $n^2$, we can conclude that by the same reasoning that $(\star\star)$ is $O(1/n^2)$. 

For the second term, note that
\begin{align*}
    K_n &= \left| \pr{\alpha^{(n)}_0}_j f^0 + \pr{\alpha^{(n)}_n}_j f^1 - \frac{\sqrt{2j-1}}{2}\pr{P_{j-1}(-1)f^0 + P_{j-1}(1)f^1} \right| \\
    &\leq M \left| \pr{\alpha^{(n)}_0}_j - \frac{\sqrt{2j-1}}{2} P_{j-1}(-1) \right| 
    + M \left| \pr{\alpha^{(n)}_n}_j - \frac{\sqrt{2j-1}}{2} P_{j-1}(1)\right|.
\end{align*}
It suffices to show that each term in the RHS is $O(1/n)$. Also note
\begin{align*}
    \pr{\sqrt{2j-1}P_{j-1}(-1)}_j &= \pr{(-1)^{j-1} \sqrt{2j-1}}_j 
    = V \text{diag} \pr{1, 0, \dots, 0} V^{-1} e_1 \\
    \pr{\sqrt{2j-1}P_{j-1}(1)}_j
    &= \pr{\sqrt{2j-1}}_j
    = V \text{diag} \pr{1, 2, \dots, N} V^{-1} e_1.
\end{align*}

% we aim to prove 
% $\abs{\pr{\alpha^{(n)}_0}_j - \frac{\sqrt{2j-1}}{2}P_{j-1}(-1) f^0}, \abs{ \pr{\alpha^{(n)}_n}_j - \frac{\sqrt{2j-1}}{2}P_{j-1}(1) f^n } = O(1/n)$.
Recall that for
\begin{align*}
        \tilde{Q}_n &:= \prod_{j=1}^n \pr{I - \frac{1}{2j}A}\\
        \tilde{R}_n &:= \prod_{j=1}^n \pr{I + \frac{1}{2j}A}^{-1}
    \end{align*}
the exact expression for the numerical solution obtained by the bilinear method is
\begin{align*}
    c^n = 
    \tilde{Q}_{n-1} \tilde{R}_{n} \pr{I+A/2} c^1 + 
    \tilde{Q}_{n-1} \tilde{R}_{n} \sum_{l=1}^{n-1} \tilde{Q}_l^{-1} \tilde{R}_l^{-1}  \frac{1}{2}\pr{\frac{1}{l}Bf^l + \frac{1}{l+1}Bf^{l+1}}.
\end{align*}
Now, consider the rightmost endpoint, i.e. the coefficient of $f^n$.
From the expression above, we can immediately find
\begin{align*}
    \pr{\alpha_n^{(n)}}_j
    &= \frac{1}{2} \pr{I+\frac{1}{2n}A}^{-1} B f^n 
    = \frac{1}{2} V \pr{I + \frac{1}{2n}D}^{-1}  V^{-1}e_1 f^n \\
    &= \frac{1}{2} V \text{diag}
    \pr{\frac{1}{1+1/2n}, \frac{2}{1+1/n}, \dots, \frac{N}{1+N/2n}}
    V^{-1}e_1 f^n.
\end{align*}
From this expression, we can directly verify that $\abs{ \pr{\alpha_n^{(n)}}_j -  \frac{\sqrt{2j-1}}{2}P_{j-1}(1)} = \mathcal{O}(1/n)$ for all $j \in \{1,2, \dots, N\}$. 

% the difference between this quantity and $\frac{\sqrt{2j-1}}{2n} f^n = \pr{\frac{1}{2n} B}_j f^n$ is $O(1/n^2)$. 
It remains to check the leftmost endpoint, i.e. coefficient of $f^0$. The terms containing $t=0$ in $c^n$ are 
\begin{align*}
    \tilde{Q}_{n-1} \tilde{R}_{n} \pr{I+A/2} c^1 
    &= \tilde{Q}_{n-1} \tilde{R}_{n} 
    \pr{ c^0 + \frac{h}{2}\pr{I+A}^{-1}Bf'(0)
    + \frac{B}{2}f^1}.
\end{align*}
Notice that due to the extra $h = \mathcal{O}(1/n)$ term,
$f'(0)$ term is negligible. Now considering the remaining term for the coefficient of $f^0$,
\begin{align*}
    \pr{\alpha_0^{(n)}}_j &= 
    n \tilde{Q}_{n-1} \tilde{R}_{n} c^0 
    = nV \pr{I+\frac{1}{2n}D}^{-1}
         \prod_{k=1}^{n-1} 
         \pr{I - \frac{1}{2k}A}\pr{I + \frac{1}{2k}A}^{-1}
    V^{-1} e_1f^0 \\
    &= nV \pr{I+\frac{1}{2n}D}^{-1}
    \pr{ \text{diag}\pr{ \prod_{k=1}^{n-1}\pr{\frac{k-1/2}{k+1/2} }, \prod_{k=1}^{n-1}\pr{\frac{k-1}{k+1} }, \dots, \prod_{k=1}^{n-1} \pr{ \frac{k-N/2}{k+N/2} }}}V^{-1} e_1 f^0 \\
    &= V \pr{I+\frac{1}{2n}D}^{-1}
    \pr{ \text{diag}\pr{ \frac{n/2}{n-1/2}, 0, \frac{(-3/2)(-1/2)(1/2)n}{(n+3/2)(n+1/2)(n-1/2)} \dots,  \prod_{k=1}^{n-1} \pr{ \frac{k-N/2}{k+N/2} }}}V^{-1} e_1 f^0.
\end{align*}
Note that the prefix term $\pr{I + \frac{1}{2n}D}^{-1}$ does not affect the asymptotic rate with respect to $n$. 
Observe that the $j$-th term in the diagonal matrix is $0$ if $j$ is an even number, and $\mathcal{O}(1/n^{j-1}$) if $j$ is an odd number. Hence  $\abs{ \pr{\alpha_0^{(n)}}_j -  \frac{\sqrt{2j-1}}{2}P_{j-1}(-1)} = \mathcal{O}(1/n)$ for all $j \in \{1,2, \dots, N\}$. 
\end{proof}

\begin{remark}
As discussed earlier, the classical technique of bounding the global error of ODEs by adding up the LTEs is not applicable due to the singularity at $t=0$. On the other hand, the quadrature formulation only requires $f$ to have bounded variation, which does not impose strong local conditions on the input function $f$.
% {\color{red} what does this mean? XXX
% while all Taylor theorem based arguments require at least the second derivative of $f$.}
% \textcolor{blue}{I was trying to emphasize that the quadrature method is better in the sense that it works for more general $f$, not only addressing the singularity at $t=0$.}
\end{remark}

\begin{remark}
    The derived convergence rates are tight in the sense that when certain polynomials are used as the input function, the global error matches the upper bound. By direct computation, one can show that the input function $f(t) = t^2$ yields a global error of $\Theta (1/n)$ for all methods except the bilinear method.
    Similarly, for the bilinear method, input function $f(t) = t^3$ attains a global error of $\Theta(1/n^2)$. 
    Notably, our rate analysis and these matching examples show that the approximate bilinear method is genuinely a first-order method while the bilinear method is a second-order method, when applied to the LegS ODE.
\end{remark}



% \begin{figure}[h]
%     \centering
%     \begin{subfigure}[t]{\textwidth}
%     \centering
%     \includegraphics[width=0.49\textwidth]{figures/fig1/smooth_function1.png}   \includegraphics[width=0.49\textwidth]{figures/fig1/smooth_function2.png}
%         \caption{Smooth input function}
%     \end{subfigure}

%     \vspace{0.5cm}

%     \begin{subfigure}[t]{\textwidth}
%     \centering
%     \includegraphics[width=0.49\textwidth]{figures/fig1/non-smooth_function1.png}   \includegraphics[width=0.49\textwidth]{figures/fig1/non-smooth_function2.png}
%         \caption{Non-smooth input function}
%     \end{subfigure}
    
%     \caption{
%     {\color{red} Change it so that there are 4 sub-plots with labels (a), (b), (c), (d)XXX}  
%     Numerical convergence behavior of the {\color{red} uniform error} of the discretization methods under various regularity properties. The numerical rates agree with theoretical estimates of Theorem~\ref{thm:conv_rate}.
%     \textbf{(a)} $f(t) = 2t^3e^{-t} $ (smooth). Bilienar exhibits $\mathcal{O}(1/n^2)$ rate while others exhibit $\mathcal{O}(1/n)$ rate.
%     \textbf{(b)} $f(t) = \frac{1}{4}\sin(10t) + \frac{1}{2} \sin(\frac{10t}{3})+ \sin(\frac{10t}{7})$ (smooth). The qualitative behavior is the same as in (a).    
%     \textbf{(c)} $f(t) = \sqrt{t}$ (non-smooth, bounded variation). All methods exhibit $\mathcal{O}(1/n)$ rate.
%     \textbf{(d)} $f(t) = t^{\frac{1}{20}} \sin(\frac{1}{t})$ (not bounded variation, Riemann integrable). All methods converge in accordance with Theorem~\ref{thm:conv_riemann}, but the rates are slower than $\Theta(1/n)$.
%     % Convergence rate of the global error of various discretization methods applied to the LegS ODE.
%     % of dimension $N=64$ on the domain $t \in [0, 2]$. \textbf{(a)} For smooth input functions, the bilinear method attains $O(1/n^2)$ convergence rate, while the other methods achieve $O(1/n)$ convergence rate. 
%     % The function plotted is \textbf{(Left)} $f(t) = 2t^3e^{-t} $ and \textbf{(Right)} $f(t) = \frac{1}{4}\sin(10t) + \frac{1}{2} \sin(\frac{10t}{3})+ \sin(\frac{10t}{7})$.
%     % \textbf{(b)} For non-smooth input functions, the usual convergence rate is not guaranteed. \textbf{(Left)} If the input function has bounded variation on the given interval, all methods acquire $O(1/n)$ convergence rate. The plotted function is $f(t) = \sqrt{t}$. 
%     % \textbf{(Right)} If the input function does not have bounded variation, there is no meaningful convergence rate. The plotted function is $f(t) = t^{\frac{1}{20}} \sin(\frac{1}{t})$ for $t>0$ and $f(0)=0$.    
%     }
%     \label{fig:fig1}
% \end{figure}



\begin{figure}[h]
    \centering

    \begin{subfigure}[t]{0.48\textwidth}
        \centering
        \includegraphics[width=\textwidth]{figures/fig1/smooth_function1.png}
        \caption{Smooth input function 1}
    \end{subfigure}
    \hfill
    \begin{subfigure}[t]{0.48\textwidth}
        \centering
        \includegraphics[width=\textwidth]{figures/fig1/smooth_function2.png}
        \caption{Smooth input function 2}
    \end{subfigure}

    \vspace{0.5cm}

    \begin{subfigure}[t]{0.48\textwidth}
        \centering
        \includegraphics[width=\textwidth]{figures/fig1/non-smooth_function1.png}
        \caption{Non-smooth input function 1}
    \end{subfigure}
    \hfill
    \begin{subfigure}[t]{0.48\textwidth}
        \centering
        \includegraphics[width=\textwidth]{figures/fig1/non-smooth_function2.png}
        \caption{Non-smooth input function 2}
    \end{subfigure}

    \caption{Numerical convergence behavior of the global error of the discretization methods under various regularity properties. The numerical rates agree with theoretical estimates of Theorem~\ref{thm:conv_rate}.
    \textbf{(a)} $f(t) = 2t^3e^{-t} $ (smooth). Bilinear exhibits $\mathcal{O}(1/n^2)$ rate while others exhibit $\mathcal{O}(1/n)$ rate.
    \textbf{(b)} $f(t) = \frac{1}{4}\sin(10t) + \frac{1}{2} \sin(\frac{10t}{3})+ \sin(\frac{10t}{7})$ (smooth). The qualitative behavior is the same as in (a).    
    \textbf{(c)} $f(t) = \sqrt{t}$ (non-smooth, bounded variation). All methods exhibit $\mathcal{O}(1/n)$ rate.
    \textbf{(d)} $f(t) = t^{\frac{1}{20}} \sin(\frac{1}{t})$ (not bounded variation, Riemann integrable). All methods converge in accordance with Theorem~\ref{thm:conv_riemann}, but the rates are slower than $\mathcal{O}(1/n)$.}
    \label{fig:fig1}
\end{figure}

% \subsection{Experimental results} \label{subsec:conv_rate_exp}
\section{Numerical experiments}
We perform simple numerical experiments to verify the tightness of the convergence theory. The results are shown in Figure~\ref{fig:fig1}. The experiments were carried out with dimension $N = 8$, a time domain $t \in [0, 2]$, and stepsize $h = 2/n$. The y-axis represents the global error $
\|c^n-c(t_n)\|$.
As discussed in the caption of Figure~\ref{fig:fig1}, the numerical behavior precisely matches the theoretical guarantees of Theorems~\ref{thm:conv_riemann} and \ref{thm:conv_rate}.
% The plot is drawn on a loglog scale to exemplify the convergence rate. Note that changing the time domain or the dimension $N$ does not qualitatively affect the result. 



% The plots in the first row show examples where the input function is sufficiently smooth, specifically $f \in C^2([0,2])$. As proven in Theorem \ref{thm:conv_rate}, the backward, approximate bilinear, zero-order hold method attains a rate of $O(1/n)$, while the bilinear method acquires the $O(1/n^2)$ rate. For non-smooth functions, however, the bilinear method is not guaranteed with $O(1/n^2)$ convergence rate. If $f$ has bounded variation, all methods acquire $O(1/n)$ rate. The left bottom plot shows the numerical error incurred for the input function $f(t) = \sqrt{t}$. Note that the nondifferentiability at the single point $t=0$ drags the convergence rate for the bilinear method. If the input function does not have bounded variation, a meaningful convergence rate is not obtainable. The bottom right plot shows the error for the input function $f(t) = t^{1/20}\sin(\frac{1}{t})$ for $t>0$ and $f(0)=0$. While the input function is continuous, one can show that it does not have a bounded variation due to the severe oscillation near $t=0$. One can see that the speed of convergence is significantly slower than $O(1/n)$, 
% without any apparent convergence rate.

% \begin{figure}[h] \label{fig:fig1}
%     \centering
%     \includegraphics[width=0.49\textwidth]{figures/fig1/smooth_function1.png}
%     % \hfill
%     \includegraphics[width=0.49\textwidth]{figures/fig1/smooth_function2.png}
%     \caption{
%     Convergence rate of the global error of various discretization methods for the LegS ODE of dimension $N=64$ on the domain $t \in [0, 2]$. (\textbf{Left}) For smooth input functions, the bilinear method attains $O(1/n^2)$ convergence rate, while the other methods achieve $O(1/n)$ convergence rate. The input function plotted is $f(t) = \frac{1}{4}\sin(10t) + \frac{1}{2} \sin(\frac{10t}{3}) + \sin(\frac{10t}{7})$.
%     (\textbf{Right}) For non-smooth integrable functions, the global error converges to $0$ as the number of timesteps $n \to \infty$, but a meaningful convergence rate is not obtainable. The input function plotted is $f(t) = t^{\frac{1}{20}}$. 
%     }
%     \label{fig:combined}
% \end{figure}


% \section{Extra Experiments (include?)}
% \textcolor{red}{Think about whether to include this part, and if so, how.}
% In this section, we experimentally verify the theoretical results proven in the paper. 

% \subsection{Pre-determined initial condition}
% In Theorem \ref{thm:exist and unique}, we have proven that the LegS ODE does not have the freedom of choosing the initial value
% the initial value at $t=0$ for the LegS ODE 




\section{Conclusion} \label{sec:conclusion}
In this work, we lay the mathematical foundations of the HiPPO-LegS ODE and its discretization methods. First, we established the existence and uniqueness of the solution of the HiPPO-LegS ODE, despite the singularity at $t=0$, and presented this result as Theorem~\ref{thm:exist and unique}. Next, we provided a framework for analyzing various discretization methods applied to the LegS ODE by reinterpreting the discretization methods as numerical quadratures. Using this insight, we proved convergence for all discretization methods of interest for any Riemann integrable input function $f$ and presented the conclusions as Theorem~\ref{thm:conv_rate}. Lastly, we established a $\mathcal{O}(1/n)$-rate under an additional bounded variation assumption on $f$ and presented the results in Theorem~\ref{thm:conv_rate}.




% This result highly depended on the structure of matrix $A$, where all eigenvalues are positive. 
% Furthermore, we provided a framework for analyzing discretization schemes, proving convergence of various discretization methods applied to the LegS ODE for Riemann integrable input functions. 


% These contributions advance the theoretical understanding of HiPPO-based models and offer a solid foundation for their continued development and application. By ensuring mathematical rigor, we pave the way for more robust and interpretable use of state-space models in machine learning, while also opening avenues for further research in both numerical analysis and the broader application of state-space representations. 


% \subsection{Figures}
% \lipsum[10]
% See Figure \ref{fig:fig1}. Here is how you add footnotes. \footnote{Sample of the first footnote.}
% \lipsum[11]

% \begin{figure}
% 	\centering
% 	\fbox{\rule[-.5cm]{4cm}{4cm} \rule[-.5cm]{4cm}{0cm}}
% 	\caption{Sample figure caption.}
% 	\label{fig:fig1}
% \end{figure}

% \subsection{Tables}
% See awesome Table~\ref{tab:table}.

% The documentation for \verb+booktabs+ (`Publication quality tables in LaTeX') is available from:
% \begin{center}
% 	\url{https://www.ctan.org/pkg/booktabs}
% \end{center}


% \begin{table}
% 	\caption{Sample table title}
% 	\centering
% 	\begin{tabular}{lll}
% 		\toprule
% 		\multicolumn{2}{c}{Part}                   \\
% 		\cmidrule(r){1-2}
% 		Name     & Description     & Size ($\mu$m) \\
% 		\midrule
% 		Dendrite & Input terminal  & $\sim$100     \\
% 		Axon     & Output terminal & $\sim$10      \\
% 		Soma     & Cell body       & up to $10^6$  \\
% 		\bottomrule
% 	\end{tabular}
% 	\label{tab:table}
% \end{table}

% \subsection{Lists}
% \begin{itemize}
% 	\item Lorem ipsum dolor sit amet
% 	\item consectetur adipiscing elit.
% 	\item Aliquam dignissim blandit est, in dictum tortor gravida eget. In ac rutrum magna.
% \end{itemize}

% \newpage
\bibliographystyle{plain}
% \bibliography{references}  %%% Uncomment this line and comment out the ``thebibliography'' section below to use the external .bib file (using bibtex) .
\bibliography{ref}


% \input{appendix}
% \include{appendix}

\end{document}
